% -*- coding: utf-8 -*-
% Preamble
\documentclass[french, 12pt]{report}
	\usepackage[utf8]{inputenc}
	\usepackage{url}
	\usepackage{etoc}
	\usepackage{float}
	\usepackage{amssymb}
	\usepackage{amsmath}
	\usepackage{fontspec}
	\usepackage{hyperref}
	\usepackage{setspace}
	\usepackage{amsfonts}
	\usepackage{graphicx}
	\usepackage{listings}
	\usepackage{bibentry}
	\usepackage{tocloft}
	\usepackage{pgfplots}
	\usepackage{titlesec}
	\usepackage{afterpage}
	\usepackage[T1]{fontenc}	
	\usepackage[round]{natbib}
	\usepackage[french]{babel}
	\usepackage[acronym]{glossaries}
	\title{Algorithmisation des modèles de confiance et performativité du système socio-fiscal : vers un paradigme disruptif pour l'appréhension des actifs numériques}
	\author{Louis Brulé Naudet\thanks{Université Paris-Dauphine}}
	\date{\today}

% Set parameters
\setcounter{page}{0}
\hypersetup{linktoc=all}

% Define style for programming displays
\lstdefinelanguage{JavaScript}{
  keywords={typeof, new, true, false, catch, function, return, null, catch, switch, var, if, in, while, do, else, case, break, const, let},
  ndkeywords={class, export, boolean, throw, implements, import, this},
  sensitive=false,
  comment=[l]{//},
  morecomment=[s]{/*}{*/},
  morestring=[b]',
  morestring=[b]"
}

\lstset{
  language=JavaScript,
  aboveskip=3mm,
  belowskip=3mm,
  showstringspaces=false,
  columns=flexible,
  numbers=none,
  breaklines=true,
  breakatwhitespace=true,
  tabsize=3,
}

\addtolength{\skip\footins}{0.6pc}
\renewcommand*\footnoterule{} %Footnode separator line

\def\siecle#1{\textsc{\romannumeral #1}\textsuperscript{e}~siècle}

\renewcommand{\cftsecleader}{\cftdotfill{\cftdotsep}} %places dots on sections lines as well

\cftsetindents{section}{0pt}{4em}
\cftsetindents{subsection}{10pt}{4em}
\cftsetindents{subsubsection}{20pt}{4em}
\cftsetindents{paragraph}{30pt}{4em}
\cftsetindents{subparagraph}{40pt}{4em}
\def\cftdotsep{1}
\cftsetpnumwidth{1em}

\renewcommand{\cftchapafterpnum}{\vspace{\cftbeforechapskip}}

\newcommand\emptypage{
	\null
	\thispagestyle{empty}
	\addtocounter{page}{-1}
	\newpage
}




\newcounter{mypara}
\setcounter{mypara}{0}

% Make Glossaries
\makeglossaries
\newacronym{Insee}{Insee}{Institut national de la statistique et des études économiques}
\newacronym{Tracfin}{Tracfin}{Traitement du renseignement et action contre les circuits financiers clandestins}
\newacronym{BTC}{BTC}{Bitcoin}
\newacronym{DEEP}{DEEP}{Dispositif d'enregistrement électronique partagé}
\newacronym{BCC}{BCC}{Bitcoin Cash}
\newacronym{PIB}{PIB}{Produit Intérieur Brut}
\newacronym{CMF}{CMF}{Code monétaire et financier}
\newacronym{BOFiP}{BOFiP}{Bulletin officiel des finances publiques}
\newacronym{ETH}{ETH}{Ether}
\newacronym{TUE}{TUE}{Traité sur l'Union Européenne}
\newacronym{TFUE}{TFUE}{Traité sur le fonctionnement de l'Union Européenne}
\newacronym{PSAN}{PSAN}{Prestataire de services sur actifs numériques}
\newacronym{AMF}{AMF}{Autorité des marchés financiers}
\newacronym{ACPR}{ACPR}{Autorité de contrôle prudentiel et de résolution}

\newglossaryentry{ipc}
{
	name=indice des prix à la consommation,
	description={Mesure statistique permettant d'estimer, entre deux périodes données, la variation moyenne des prix des produits et des services consommés par les ménages, à qualité constante.
	}
}
\newglossaryentry{désutilité marginale du travail}{
	name=désutilité marginale du travail,
	description={Insatisfaction supplémentaire en termes d'utilité, apportée par l'acquisition d'une unité de travail.
	}
}
\newglossaryentry{digeste}{
	name=Digeste,
	description={[\emph{Nom propre}] Recueil méthodique d'extraits des opinions et sentences des juristes romains, réunis sur l'ordre de l'empereur Justinien en 1530 (publié le 16 décembre 1533). Fondamental dans la compréhension historique des idées juridiques, le Digeste forme la deuxième partie du \emph{Corpus iuris civilis}, et se composera de 9000 fragments rédigés par 40 juristes, affirmant certains aspects de la procédure archaïque, du droit des choses, des actions et des obligations, mais également de la possesion, de l'\emph{animus}, et de la distinction entre le droit écrit et non écrit.
	}
}
\newglossaryentry{courbe d'indifférence}{
	name=courbe d'indifférence,
	description={Une courbe d'indifférence est le lieu géométrique (une courbe dans le plan) des paniers entre lesquels l'individu est exactement indifférent. La courbe d'indifférence associée à $x$ est définie comme l'ensemble des paniers y donnant le même niveau de satisfaction que le panier x :
$ \lbrace y/y \backsim x \rbrace $.
	}
}
\newglossaryentry{hard fork}{
	name=hard fork,
	description={En bases de données décentralisées, la notion de \emph{hard fork} consiste en un changement du paradigme de validation d'ensembles de transactions, partagé par tous les noeuds du réseau, pouvant conduire à une division permanente de la chaîne de blocs en cas de concurrence entre différentes versions logicielles.
	}
}
\newglossaryentry{convertibilité}{
	name=convertibilité,
	description={Se définie comme étant la propriété que possède une monnaie nationale d'être librement échangée, la convertibilité s'opère, selon les époques, soit contre de l'or dans le pays d'émission (ou de l'argent dans certains pays), soit contre des monnaies étrangères.
	}
}
\newglossaryentry{taux d'escompte}{
	name=taux d'escompte,
	description={Le taux d'escompte est un taux d'intérêt utilisé sur le marché monétaire, pour les prêts à très court terme (quelques jours).
	}
}
\newglossaryentry{thésaurisation}{
	name=thésaurisation,
	description={Comportement économique résultant de la volonté de conserver la valeur (monnaie, billets, pierres précieuses, or...) en dehors du système économique pour en tirer un profit ou par absence de meilleur emploi.
	}
}
\newglossaryentry{Loi de Gresham}{
	name=Loi de Gresham appliquée au bimétallisme or-argent,
	description={La plus importante rareté relative de l'argent par rapport à l'or, suite à la découverte d'importants gisements en Californie et en Australie dès 1848, conduit à un dépassement de la valeur légale de l'argent par sa valeur physique. Les agents économiques se retrouvent dans une situation ou il devient préférable d'acheter de l'argent avec de l'or au cours légal, afin de le revendre sur les marchés, la banque centrale réglant la différence de valeur par la liquidation de ses réserves d'argent. La reconstitution des réserve entraine un mécanisme plus complexe : la banque centrale se trouve en une situation contraignante, et l'acquisition impérieuse d'argent au prix des marchés conduit les agents à spéculer sur sa réévaluation. Le gain à l'échange couvrant les coûts de refonte, de frappe et de transport, la demande de thésaurisation des monnaies d'argent entrainera à moyen terme une crise des règlements par la diminution relative de circulation des signes monétaires.
	}
}
\newglossaryentry{avilissement}{
	name=avilissement,
	description={Action de se déprécier, de perdre de sa valeur; résultat de cette action.
	}
}
\newglossaryentry{double dépense}{
	name=double dépense,
	description={Le problème de double dépense renvoie, en informatique des systèmes décentralisés, à une situation ou un même débiteur serait en position de signer en envoyer sans intervalle de temps, deux transactions prenant pour point d'entrée une unique unité de monnaie-cryptographique, à destination de deux bénéficiaires distincts. En conséquence, la probabilité pour que deux ensembles de \emph{n} n\oe{}uds connaissent d'une information divergente sur l'état de la chaîne de transactions pendante accentue le risque de désaccord sur la validité du registre à l'échelle du réseau.
	}
}
\newglossaryentry{système pair à pair}{
	name=système pair à pair,
	description={Décrit un modèle d'échange en réseau où chaque entité est à la fois client et serveur. Dans le cadre du Bitcoin, ce système peut être qualifié de parfaitement décentralisé du fait que les connexions entre participants se réalisent sans intervention d'une infrastructure de répartition (tiers de confiance). Le partage d'information, du fait de la dynamique réticulaire, gagne en performance et en immutabilité : une censure ou tentative de modification des transaction ne sera efficiente qu'à condition de suspendre l'interconnexion entre les n\oe{}uds.
	}
}
\newglossaryentry{maximum supply}{
	name=maximum supply,
	description={Indicateur économique de la quantité maximale théorique d'émission d'un jeton de monnaie-cryptographique.
	}
}
\newglossaryentry{halving}{
	name=halving,
	description={Phénomène périodique de réduction de moitié de la prime de résultat associée au processus de validation d'un bloc de transactions du DEEP. Le halving s'interprète en un mécanisme de convergence de la masse monétaire vers l'offre maximale théorique imprimée lors du développement du protocole.
	}
}
\newglossaryentry{Livre blanc}{
	name=Livre blanc,
	description={Document d'information concis traitant des objectifs et moyens mis en \oe{}uvre à la production d'un dispositif d'enregistrement électronique partagé. Dans le cadre du règlement MiCA, ce dernier devant notamment contenir une description détaillée de l'émetteur, des droits, obligations et risques attachés aux crypto-actifs, ainsi qu'une présentation des principaux participants à la conception et à la mise au point du projet.
	}
}

\setcounter{secnumdepth}{4}

% Document
\begin{document}{
\sloppy
% Titlepage
\begin{titlepage}
\begin{center}
{Université Paris-Dauphine (Paris-Sciences et Lettres)} \vspace{1.5 cm}\\
\end{center}
\begin{center}
\Huge{La qualification juridique des monnaies-cryptographiques au miroir de la concurrence des devises}
\end{center}
\vspace{0cm}
\begin{center}
\Large{De l'influence des usages et du régime fiscal sur l'évaluation du caractère monétaire}
\end{center}
\vspace{1.5 cm}
\begin{center}
\normalsize{Louis Brulé Naudet \\ Sous la direction de Morgan Sweeney}
\vspace{1.5 cm}
\end{center}
\begin{center}
Centre de Recherche Droit Dauphine (CR2D)\\
Place du Maréchal de Lattre de Tassigny\\
75016 Paris, France
\end{center}
\vspace{1.5 cm}
\begin{center}
Mémoire de recherche présenté au département de sciences des organisations en vue de l'obtention du grade Master en Droit des affaires et fiscalité, programme gradué Droit.
\end{center}
\vspace{1.5 cm}
\end{titlepage}

\emptypage
\emptypage
\nobibliography*
\setcounter{tocdepth}{4}
\setcounter{secnumdepth}{4}

\titleformat{\chapter}
 {\normalfont\fontsize{21}{23}\mdseries}{\thechapter}{1em}{}
\titleformat{\section}
  {\normalfont\fontsize{16}{19}\mdseries}{\thesection}{1em}{}
\titleformat{\subsection}
  {\normalfont\fontsize{14}{17}\mdseries}{\thesubsection}{1em}{}
\titleformat{\subsubsection}
  {\normalfont\fontsize{13}{16}\mdseries}{\thesubsubsection}{1em}{}
\titleformat{\paragraph}
 {\normalfont\fontsize{12}{15}\mdseries}{\theparagraph}{1em}{}
\titleformat{\subparagraph}
 {\normalfont\fontsize{11}{14}\mdseries}{\thesubparagraph}{1em}{}

% Table of contents
\tableofcontents
\chapter{Introduction}
\og Pour dire encore un mot sur la prétention d'enseigner comment doit être le monde, nous remarquons qu'en tout cas, la philosophie vient toujours trop tard. En tant que pensée du monde, elle apparaît seulement lorsque la réalité a accompli et terminé son processus de formation. Ce que le concept enseigne, l'histoire le montre avec la même nécessité : c'est dans la maturité des êtres que l'idéal apparaît en face du réel et après avoir saisi le même monde dans sa substance, le reconstruit dans la forme d'un empire d'idées. Lorsque la philosophie peint sa grisaille dans la grisaille, une manifestation de la vie achève de vieillir. On ne peut pas la rajeunir avec du gris sur du gris, mais seulement la connaître. Ce n'est qu'au début du crépuscule que la chouette de Minerve prend son vol.\fg\par

\textemdash Georg Wilhelm Friedrich Hegel, \emph{Principes de la philosophie du droit}, Gallimard, 1940.\par
% First chapter
\chapter{Une tendance à l'érosion des différences fonctionnelles entre monnaie légale et monnaie cryptographique : vers un commerce juridique plurimonétaire en devenir ?}
\flushbottom 
\setstretch{1.42}
\og Si l'usage a créé des mots spéciaux, ou bien s'il a été détourné de leur sens primitif des mots qui à l'origine exprimaient des poids, c'est que l'unité est arrivée à représenter autre chose qu'une certaine quantité de métal, et qu'à cette chose nouvelle, il a fallu un nom nouveau.\fg\par
\textemdash Marcel Mongin, Professeur à la faculté de Droit de Dijon, \emph{Des changements de valeur de la monnaie}, 1887.\par

\setstretch{1}\section{Réflexion sur les attributs fonctionnels de la monnaie}\setstretch{1.42}

La monnaie, qu'importe ses qualités structurelles et au regard de tout l'ensemble des choses du commerce, possède une valeur stable à un moment donné \citep{mongin1887}. En guise de propos liminaires à notre recherche, il semble fondamental d'éclaircir certaines propositions pourtant élevées en qualité d'axiomes, concernant les attributs fonctionnels de l'objet monétaire.\par

\setstretch{1}\subsection{Isomorphisme d'ordre et conception théorique d'une monnaie de compte}\setstretch{1.42}

Lorsque l'on distend les liens constitutifs de la chose monétaire, l'anthropologie économique consacre avec une particulière attention la problématique de quantification\footnote{\bibentry{dosquet2018}.}, en tant qu'ouverture théorique vers le dilemme de l'articulation et de l'objectivation de la valeur. Penser la fixité de la valeur de la monnaie permet de sacraliser l'homomorphisme bijectif, comme correspondance justifiant de conserver la relation combinatoire entre le prix et la valeur de la chose. Alors, il se présente essentiel de l'expliciter, afin de comprendre les causes de l'inflation soulevées par l'orthodoxie, et d'apprécier les origines de la théorisation d'une monnaie de compte.\par
	
\setstretch{1}\subsubsection{Neutralité économique de la monnaie et distorsion des prix}\setstretch{1.42}

Bien qu'à première vue, il apparaisse paradoxal de considérer que les variations de prix ne soient pas dues aux jeux de la dépréciation et de son contraire, en conformité avec la réalité économique, la théorie générale sur la fixité de la valeur présente des qualités indéniables pour la représentation des faits contemporains. Dans un monde de marchés \footnote{\bibentry{guesnerie2006}.}, composé de l'ensemble des biens meubles et immeubles, ainsi que de l'universalité des contrats de présentations de services, l'idée d'un changement de valeur de la monnaie devrait, de manière spontanée et parfaitement rationnelle, exercer son inspiration sur tous les prix sans exception, et ce, quel que soit l'objet ou la nature de l'échange. Dans une logique quasiment physique, prenant pour objet d'étude la dynamique des prix, le propre même de cette aspiration devrait se matérialiser dans une unicité de direction et de sens, avec pour réflexion, une égalité des taux de croissance de l'ensemble considéré. Cette dynamique est-elle ce que nous offre la réalité ?\par

Attendu que l'on s'essaye à une analyse pratique de l'évolution des prix, tout au contraire, on remarque aisément que les variations s'opèrent, d'une part, dans des proportions inégales, mais surtout, dans des directions parfois divergentes. Si l'on analyse les différents groupements conjoncturels donnés par les indices des prix à la consommation de l'\acrshort{Insee} pour le mois d'août 2021\footnote{\gls{ipc}, résultats définitifs (IPC), Août 2021, Numéro 234.}, on observe parallèlement des mouvements opposés d'une forte intensité, ne trouvant pas, par définition de l'indicateur statistique, leur source dans la qualité et la quantité des produits considérés. Par ces constats en apparence accessibles à tout individu participant à la vie économique de la cité, se met naturellement en exergue, l'une des controverses majeures de la macroéconomie moderne. Comment expliquer l'absence de cohérence apparente dans les variations catégorielles des prix ? Comment est-il imaginable, que des contractants, indifférents, en pratique, aux statistiques de la circulation monétaire, subissent contre leurs desseins, une modification des prix ?\footnote{\bibentry{friedman1959}.} Une des réponses qui semble pertinente pour l'explication de ce phénomène consisterait à minorer drastiquement les effets dûs à une dépréciation ou à une hausse de la monnaie, au profit d'un changement de valeur qu'éprouveraient les objects du commerce. Si l'on constate qu'une unité marchande de blé ou qu'une heure de travail possèdent des prix qui diffèrent d'une période à une autre, il est concevable de considérer que des causes spéciales et relatives à ces objets ont participé à une modification de leurs valeurs à l'égard de l'ensemble des autres choses du commerce. En conséquence, plusieurs facteurs explicatifs peuvent être soulignés dans l'interprétation des mouvements de prix. Dans le cadre de notre propos introductif, nous en retiendrons deux, l'effet de la demande, et l'effet des ajustements conjoncturels sur la structure de coûts des entreprises. Dans un sytème micro-fondé, il est en réalité relativement aisé de comprendre l'origine non monétaire de l'inflation. En 1958, Alban William Phillips\footnote{\bibentry{phillips1958}.} établira, en ce sens, une relation empiriquement négative entre la croissance des salaires et le taux de chômage au Royaume-Unis, prenant la forme suivante :\par

\[  \frac{\Delta W}{W}= \alpha.U^{-\beta}- \gamma \]
	
Dans laquelle \emph{U} représente le taux de chômage, et le membre de droite de l'équation, le taux de croissance du salaire nominal annuel. Sans mathématiser davantage le problème, cette équation simplifiée nous permet de fixer un mécanisme économique complet, fondé sur l'existence de pressions sur le marché du travail\footnote{\bibentry{samuelson1960}.}. Ainsi, elle permet de démontrer qu'une hausse du niveau d'emploi engendre tendanciellement un accroissement de la \gls{désutilité marginale du travail} et conduit les syndicats à prétendre à des augmentations de rémunérations, accélérant mécaniquement le coût marginal de production et conduisant les producteurs en concurrence monopolistique à élever le prix de leurs marchandises. En parallèle, la hausse de la consommation autorisée par la croissance des salaires exerce un effet similaire, rendant, selon les propensions individuelles, soit le travail plus coûteux en termes d'utilité, soit le loisir plus attractif. Sans prendre nécessairement en compte l'existence d'une fiscalité distorsive reportée sur les salaires par les associations syndicales, nous venons d'illustrer les deux causes d'inflation non monétaire citées en amont de notre propos. D'une part, l'augmentation de la rémunération induit une modification structurelle du panier de consommation des travailleurs, pesant sur le jeu de l'offre et de la demande avec l'émergence de biens polarisant les désirs individuels. D'autre part, la modification de la structure des coûts des entreprises se répercute sur le prix des choses du commerce, des entreprises $\psi$, révisant leurs grilles tarifaires de manière optimale. Certes, ce modèle simplifié ne représente pas toutes les interactions de l'équation de détermination du niveau général des prix, toutefois, il permet de porter un regard pluridisciplinaire sur des problématiques croisées, et démontre combien est fragile la théorie générale sur l'instabilité de la monnaie.\par
	
\setstretch{1}\subsubsection{Du modèle ordinal à l'émergence d'une unité de compte abstraite}\setstretch{1.42}

En réalité, il nous reste un fait, en apparence, bien ésotérique à expliquer, que celui d'un objet conservant sa valeur malgré l'incohésion des oscillations que subissent tous les autres. Poussons le questionnement encore plus profondément, comment expliquer que l'or et l'argent, deux matériaux consacrés par le système bimétallique\footnote{Loi du 7-17 germinal an \textsc{xi} (28 mars-7 avril 1803), consacrant le franc germinal.} du \siecle{19}, exercent à la fois le rôle de biens désirables soumis aux lois de l'offre et de la demande \footnote{\bibentry{marshall1890}.}, et celui d'instruments monétaires dont nous avançons pouvoir justifier  la stabilité ? Voilà une marchandise qui éprouverait des changements de valeur suivant la forme sous laquelle elle circule, en se présentant successivement comme fixe lorsque sa fonction n'est pas dévolue aux usages monétaires, et sensible aux fluctuations dès lors qu'elle présente une impossibilité de conversion en monnaie au gré des agents économiques qui la possèdent. On commence à percevoir l'acceptabilité de conception d'une abstraction parfaite entre l'unité de compte et le médium matériel ou la structure ne semble en rien être à l'origine d'une mutation de la valeur\footnote{Dans notre exemple, de la monnaie bimétallique.}. Replaçons-nous dans notre environnement micro-fondé en qualité d'agent économique dont le revenu est égal à la valeur de marché de sa dotation\footnote{On supposera que la dotation de l'agent est constituée de sa rémunération du travail \(p_1w_1\) ainsi que d'un patrimoine initial parfaitement exogène \(p_nw_n\), tous deux exprimés en termes réels. Celui-ci ne contrôle aucunement sa dotation, et n'a pas la possibilité d'exercer son influence, dans l'hypothèse de marchés purement concurrentiels, sur le niveau des prix.} $\delta$ telle que :
	
\[ \delta = p_1w_1 + ... + p_nw_n\]

Si l'on suppose que l'environnement institutionnel de l'agent se limite, pour l'essentiel, au marché concurrentiel\footnote{Est ainsi faite abstraction du poids d'une quelconque entité étatique.}, et que les prix (\(p_1;..;p_n\)) conservent une parfaite indépendance à l'égard des quantités de biens consommées (\(x_1 + ... + x_n\)), alors, la valeur de marché du plan de consommation économiquement réalisable\footnote{Nous ne démontrons ici que la formule générale, toutefois, on notera qu'il existe une infinité de plans de consommation économiquement réalisables.} s'obtiendrait par :

\[ p_1x_1 + ... + p_nx_n\leq  \delta \]
 
 Persiste toutefois un élément de généralisation qui doit être résolu. La question de la détermination du prix permet de soulever une multitude de contraintes pesant sur la désirabilité relative des consommations alternatives. En ce sens, il est par exemple impossible pour l'individu de consommer une valeur de marché strictement supérieure au plan de consommation égalisant celle de sa dotation. Considérons désormais une valeur de \emph{n} suffisamment grande pour couvrir la constellation des marchandises du commerce. Dans les fondements théoriques de notre analyse, nous avons omis volontairement l'existence d'une monnaie légale comme instrument de paiement, afin que le transfert des ressources économiques ne puisse se produire que par l'échange de choses du commerce, désirables pour leurs qualités intrinsèques\footnote{\bibentry{say1803}.}. Tous sont ainsi contraints d'effectuer un jeu de possessions et de dépossessions successives afin de parfaire l'allocation de chacun, considérée comme durablement inégale en tant que telle, du fait de l'absence d'intervention publique. Empiriquement, il sera possible de procéder à une évaluation quantitative des flux relatifs à chaque marchandise, de sorte que naisse une relation de classement agrégée des préférences. Ces choses seront entre elles comme les chiffres 1, 6, 7, 8... Et la détermination de la valeur d'une marchandise nouvelle découlera de sa comparaison avec celles qui ont auparavant été vérifiées. Ainsi, la grande bibliothèque des prix sera continuellement alimentée en nouveaux chiffres, qui pourront, à leur tour, servir de points de comparaison. Nous venons de décrire un modèle convergent fondé sur un univers ordinal indiquant parfaitement les rapports de valeurs. Si maintenant, pour l'accessibilité du langage, est opéré un changement de variable introduisant un nouvel ensemble, permettant d'indexer à ces chiffres un nom quelconque, les relations d'ordre entre les valeurs seront parfaitement conservées ; on assistera à la création d'une unité de compte abstraite et strictement immatérielle, ne reposant à priori sur aucune structure physique. Plutôt que de dire que les marchandises \emph{a}, \emph{b} ou \emph{c}, valent un certain rang de l'ensemble ordinal, exprimé sous la forme d'un simple chiffre, sera apposé à ce chiffre la notion d'une entité particulière, celle d'\emph{euro} ou tout autre.\par
 
 \[ \begin{matrix}
	a & b & c & d & e & \cdots & n \\
	1 & 2 & 3 & 4 & 5 &\cdots & \omega & euro
	\end{matrix} \]
	
Voilà une démonstration qui met en lumière l'une des propriétés les plus remarquables des instruments monétaires : l'unité de compte se dégage spontanément de l'ensemble des valeurs \footnote{M. Cournot, \emph{Principes de la théorie des richesses}, 1863.}. On retrouvait déjà les prémices de cette théorie dans l'\oe{}vre de Jean-Baptiste Say, pourtant, il est surprenant de constater qu'elle reste peu connue du grand public, alors même qu'elle représente le principal fondement conceptuel des monnaies modernes.\par
 
 \setstretch{1}\subsection{Approche symbolique de la monnaie et moyen d'échange des valeurs}\setstretch{1.42}

En son caractère itératif, le procédé d'évaluation comparative des différents objets du commerce imprime aux instruments monétaires une nécessité, celle de représenter symboliquement un critère d'unité au sein d'une totalité sociale\footnote{\bibentry{theret2007}.}. Le consensus autour d'un invariant\footnote{Relatif si l'on fige les effets de l'inflation et du commerce international.} structurant, pouvant aisément être interprété comme rapport d'équivalence entre les choses, permet alors de considérer la monnaie comme un vecteur efficace des dettes et des créances au sein des différents noeuds d'une société. La problématique de qualification des objets monétaires a posé d'importants questionnements dans les histoires juridiques et économiques.\par

\setstretch{1}\subsubsection{La monnaie, le bien, la chose, une approche économique à l'épreuve du Droit}\setstretch{1.42}

Alors que le Doyen Carbonnier proposait que \og ce qui confère à une chose le caractère juridique de monnaie, c'est qu'elle est reçue dans les paiements non pas pour ce qu'elle représente matériellement, mais en tant qu'équivalent, fraction ou multiple d'une unité idéale\fg{}\footnote{\bibentry{carbonnier2004}.}, il semblerait possible de transcender cette approche définitionnelle en s'appliquant à une étude interdisciplinaire de la désirabilité relative des biens, au regard de l'analyse des idées juridiques issues du \gls{digeste},  et de la théorie micro-économique du consommateur. Au sein de ce développement, nous nous attacherons à une compréhension de la notion de bien, de celle de la chose et de la translation entre le domaine juridique et économique, afin de démontrer que la monnaie est une construction exceptionnellement complexe, dotée de propriétés uniques par opposition avec la constellation des biens du commerce. La doctrine juridique s'entend de façon satisfaisante autour de la distinction entre la chose et le bien, et plus précisément, sur l'inclusion du bien au sein de l'espace des choses\footnote{\bibentry{atias2004}.}. Au sens philosophique, on admettra volontiers que se dit \emph{chose}, l'abstraction dans son acception la plus large, indifférente de tout, et dont la signification découle de la matière dont on traite. Le droit, dans son approche réaliste, apporte quelques précisions supplémentaires, en laissant entendre que la chose serait davantage le qualificatif de l'objet conceptuel présent dans la nature, et possédant une quelque utilité aux hommes, peu important son caractère appropriable duquel émanerait la possession\footnote{Vinnius, \emph{De rerum divis}. Heineccius, \emph{Elementa juris}, \no 312.}. Penser l'articulation entre la \emph{chose} et le \emph{bien} consiste à affronter le problème de la propriété, en ce que la signification de ce dernier est bien moins étendue et se matérialiserait, pour les jurisconsultes du moins\footnote{étymologiquement, le mot \emph{bien} trouverait ses racines dans le terme latin \emph{beare} infinitif de l'expression \og rendre heureux \fg{}.}, par l'introduction d'un facteur d'accroissement du bonheur de la vie, en procurant un moyen factuel d'augmenter et de conserver les jouissances\footnote{\bibentry{doss1842}.}. étrangement, le législateur semble employer indifféremment les deux termes\footnote{Voir les articles 544, 546, 548 et suivants du Code civil.}, si bien qu'il est possible de considérer les biens sous deux rapports principaux\footnote{\bibentry{duranton1844}.}, ceux qui consistent en un droit, et les choses ordinaires, intrinsèquement douées de la \emph{proprietas}\footnote{Définition présente au sein du Digeste : 22, 1, 38, 2. Par opposition à l'usufruit : \emph{Institutes} de Gaïus, 2, 33.}. La confusion entre les notions de droit atteste de la difficulté de qualification des objets qui nous entourent, si bien que l'on en vient régulièrement à distendre les effets de ces terminologies en considérant la \emph{proprietas} comme un droit\footnote{\bibentry{gaffiot2000}.} alors qu'elle serait davantage, dans sa genèse lexicographique, une conception objective pour désigner la composante de la chose appropriable\footnote{\bibentry{zenati2006}. Le droit romain, attachait en ce sens une spécificité unique à la \emph{proprietas}, la rendant insusceptible de se transformer en un droit, de sorte à ce qu'à l'occasion d'un démembrement, le nu-propriétaire ne conservait que le domaine de sa propriété.}. Une série de questions peut alors nous surgir : la monnaie est-elle un bien, une chose ou une virtualité plus complexe ? Si tel est le cas, les monnaies cryptographiques partagent-elles les propriétés les plus particulières des monnaies historiques ? Première observation, la somme d'argent, en son caractère de \emph{quantitas}, existe indépendamment des instruments monétaires dont elle fait l'addition, ou de la créance qui lui sert de véhicule dans le commerce juridique\footnote{Ibid, p. 1549. \og Il est des cas où son existence se détache très clairement : ainsi, lorsqu'un testateur fait un legs particulier de somme d'argent à prélever sur l'actif successoral\fg{}.}. De cette abstraction conceptuelle découle un principe essentiel : le paiement au sens juridique possède une étendue bien plus restreinte qu'en science économique, par le fait qu'il implique dans son essence, l'existence préalable d'un lien obligatoire\footnote{\bibentry{pignarre2010}.}. Cela se vérifie au prisme de l'herméneutique, à la lettre de l'article 1343-3 du Code civil, et ne présente pas d'équivoque, la monnaie n'est alors qu'un moyen de paiement parmi d'autres permettant d'éteindre une obligation, et lorsqu'elle est envisagée en tant que telle, l'étude fonctionnelle s'opacifie du fait que la terminologie utilisée par le législateur conduit à la confondre avec les outils qui permettent son utilisation\footnote{Article L.311-3 : \og Sont considérés comme moyens de paiement tous les instruments qui permettent à toute personne de transférer des fonds, quel que soit le support ou le procédé technique utilisé\fg{}.}. Deuxième observation, en analyse des obligations monétaires, la réalité économique détermine parfois la réalité juridique, et sa prégnance se retrouve lorsque l'on s'attache à une qualification de la nature des paiements en monnaie scripturale depuis les débuts du \siecle{20}\footnote{\bibentry{lyon1923}. \bibentry{drouillat1931}.}. L'idée que le titulaire d'un compte bancaire détient une créance sur l'établissement de crédit ne soulève guère d'interrogation particulière tant cette dernière semble transparaitre à la lecture de l'article L.312-1-1 du Code monétaire et financier\footnote{Article L.312-1-1 : \og Lorsque l'établissement de crédit est amené à proposer à son client de nouvelles prestations de services de paiement dont il n'était pas fait mention dans la convention de compte de dépôt[...]\fg{}.}. Pourtant, le paiement par virement bancaire ne saurait être interprété en une cession de créance, et la Cour de cassation, reconnaissant que l'inscription en compte, parce qu'elle correspond à la mise à disposition d'une somme d'argent au profit du titulaire, vaut paiement\footnote{Cass. 1re civ., 23 juin 1993.}, permet d'observer un processus remarquable : il ne fait aucun doute que le solde disponible sur un compte bancaire s'interprète en une créance, toutefois, l'objet de droit, et on le sent fortement, bien que confusément, a subi une véritable mutation afin de s'adapter au rôle de monnaie qui lui est dévolu de fait\footnote{J.-L. Rives-Lange, M. Cabrillac, Rép. Com., \og Virement \fg{}, \no 24 et 25. Cité par L.-F. Pignarre, op. cit., p. 56.}. En d'autres termes, on observe la naissance d'une monnaie scripturale chiffrée et non frappée, d'une fongibilité parfaite\footnote{Cette caractéristique propre à la monnaie résiderait dans sa faculté d'être remise immédiatement en circulation au sein du commerce juridique.}, incorporée dans le symbole qui en imprime le pouvoir de monstration et en révèle l'existence : l'inscription en compte. L'appui de la micro-économie pour l'effort de qualification relève une autre caractéristique quasiment exclusive à la monnaie\footnote{Considération prise que la monnaie ne doit en aucun cas être considérée comme une invention des sociétés marchandes modernes. \bibentry{theret2008}.}, et à cet effet, il convient de serrer de plus près une approche centrée sur la \gls{courbe d'indifférence}\footnote{\bibentry{edgeworth1881}. Une courbe d'indifférence est le lieu géométrique (une courbe dans le plan) des paniers entre lesquels l'individu est exactement indifférent. La courbe d'indifférence associée à $x$ est définie comme l'ensemble des paniers y donnant le même niveau de satisfaction que le panier $x$ :$\lbrace y/y \backsim x \rbrace$.} pouvant être attaché à sa détention, et à laquelle on aurait soustrait les effets du consensus au sein d'un groupe. De cette abstraction théorique et au regard de notre propos antérieur, cela nous conduirait à la possession d'un objet marquant par sa virtualité et doté d'une fonction d'évaluation n'admettant aucun support de comparaison reconnu par les autres membres : la monnaie ne permettant plus l'extinction de dettes, se voit imprimée d'une neutralité\footnote{\bibentry{say1803}.} parfaite rendant sa désirabilité relative nulle. Dans le prolongement de cette idée, si l'on décide de matérialiser notre monnaie reconnue d'une seule personne, en la confondant avec un objet qui \emph{ceteris paribus}, n'admet aucune source de bien-être, la recherche de sa détention\footnote{En posant l'hypothèse que la détention induit un appauvrissement des ressources, ou \emph{a minima} une consommation d'unités de temps pesant sur le consommateur et son cycle de vie.} pourrait conduire à un processus surprenant, celui de l'opposition du coefficient directeur de la courbe d'indifférence rattachant la monnaie à l'ensemble des biens du commerce juridique désirables $\delta$. Alors que la décroissance de la courbe d'indifférence semble être élevée au rang d'axiome\footnote{En raison des hypothèses de non saturation des préférences et de complétude, le consommateur désireux de maintenir son niveau d'utilité doit substituer un bien à un autre au sein de son panier de consommation.}, formuler mathématiquement et de façon rationnelle son antagoniste consisterait à admettre qu'il existe une infime proportion d'objets\footnote{\bibentry{delaplace2017}. La qualification de \emph{bien} pourrait être remise en question en ce sens.} dont la consommation induit une diminution non équivoque de l'utilité. Réciproquement, une baisse de la consommation de cet objet se présenterait comme un facteur d'accroissement de la satisfaction et on en déduit que le long de la courbe d'indifférence, le sens de variation de l'objet \og indésirable \fg{} est identique à celui du bien désirable : la courbe d'indifférence est ainsi croissante \footnote{$y$ est une fonction croissante de $x$ si $\Delta x \gtrless{} 0 \Rightarrow \Delta y \gtrless{} 0$}.\par

\pgfplotsset{compat=newest}
\usepgfplotslibrary{fillbetween}
\hspace{1cm}
\begin{center}
\begin{tikzpicture}
\begin{axis}[
    axis lines = middle,
    xlabel = {$monnaie$},
    ylabel = {$\delta$},
    xmin=0, xmax=5.5,
    ymin=0, ymax=24,
   	yticklabels={,,},
   	xticklabels={,,}]
 
\addplot [name path = A,
    -latex,
    domain = 0:4.5,
    samples = 1000] {x^2} 
    node [very near end, right]{};
\end{axis}
\end{tikzpicture}
\\Figure 1 \---{} Courbe d'indifférence associée à la conception d'une monnaie privée de consensus.
\end{center}

Lorsque nous affrontons à nouveau la question de la nature fondamentale de la monnaie, les évidences juridiques tendent indéniablement à s'effacer, emportées par la subversion d'un objet qui, privé de reconnaissance collective, dérogerait à l'\emph{espace fini} de la conception des choses du droit romain.\par
 
\setstretch{1}\paragraph{La monnaie, phénomène anthropologique universel}\setstretch{1.42}

La prise en considération des éléments précédents conduit au constat que la monnaie suppose une acceptation généralisée des individus qui font société. Au travers de cette conception, se fait jour une véritable approche symbolique de la monnaie dont la spécificité résiderait précisément dans le fait qu'elle cristalliserait une suspension des facultés critiques individuelles\footnote{\bibentry{orlean1991}.} : l'adhésion par médiation à la règle monétaire procèderait d'un assujettissement mécanique plus que d'un calcul rationnel. On prend là conscience du rôle de l'extériorité sociale dans l'écriture de l'épaisseur conceptuelle attachée à la monnaie, et un paradoxe reposant sur une confusion du raisonnement déductif\footnote{\bibentry{russell1908}.} semble se dessiner : un objet monétaire est accepté en paiement, à la seule condition que le tout, entendu comme l'ensemble des échangistes, l'accepte également\footnote{\emph{Ibid}., p.129.}. L'effet emporte un mécanisme cognitif qui déjà engage la cause, on retrouve la notion d'\emph{anticipation infinie des chemins parcourables}\footnote{\bibentry{merot1999}.}, et la monnaie, interprétée comme un \emph{phénomène} social, au même titre que le langage, dont elle serait une forme particulière, ne saurait se conformer à la définition traditionnelle de la \emph{chose}. Cette lecture schopenhauérienne\footnote{\bibentry{bouriau2002}.} du caractère trouble dans le rapport entre la cause et l'effet semble porter l'intégralité des questionnements autour de la qualification de l'objet monétaire, en ce sens qu'elle décrirait le paradoxe comme l'imperfection laissée dans le monde pour attester qu'il n'est rien d'autre qu'une invention de l'homme\footnote{\bibentry{borges1957}. \og C'est nous qui avons rêvé l'univers. Nous l'avons rêvé solide, mystérieux, visible, omniprésent dans l'espace et fixe dans le temps ; mais nous avons permis qu'il y eût à jamais dans son architecture de minces interstices de déraison, pour attester de sa fausseté \fg{}.}\par 
 
De cette démonstration complexe, on entrevoie analogiquement un raisonnement qui permettrait de s'approcher efficacement d'une grille de qualification systématique pour l'analyse des objets à tendance monétaire. Pourtant, la réalité juridique semble encore latente malgré les efforts du législateurs pour clarifier le régime applicable à une forme particulière de ces objets à vocation monétaire : les \emph{monnaies-cryptographiques}.\par
 
\setstretch{1}\section{Un objet juridique hors des définitions traditionnelles du droit}\setstretch{1.42}
 
Les monnaies cryptographiques ne sont plus une innovation de court terme  dans la sphère d'influence du droit, toutefois, leur appréhension par le législateur et l'administration fiscale est le sujet de nombreuses interrogations qu'il conviendra d'étudier dans le cadre de notre propos introductif, en débutant notre analyse par une brève présentation de la chronologie moderne des mutations organiques de la monnaie.\par
 
\setstretch{1}\subsection{Comprendre l'origine des mutations organiques de la monnaie dans une approche historique}\setstretch{1.42}

Difficile de ne pas se passionner pour l'histoire monétaire occidentale tellement elle semble riche et contradictoire. Nous l'avons brièvement évoquée, la stabilité monétaire connait un tournant majeur avec la promulgation de la loi du 7 germinal an \textsc{xi}\footnote{Bien que le franc était considéré comme l'unité monétaire légale de la France depuis le 28 thermidor an \textsc{iii}, au même titre que se trouvait sacralisée la division en dix décimes de dix centimes chacun, il faudra pourtant attendre le Consulat pour connaître de l'émission de la première pièce issue du monnayage du franc. \bibentry{desrousseaux2012}.} instituant le franc germinal et la consécration de l'argent comme étalon\footnote{L'or lui étant subordonné. \og Cinq grammes d'argent au titre de 9/10e de fin, constituent l'unité monétaire qui consacre le nom de franc \fg{}.}. En assignant à la monnaie un poids fixe de métal précieux\footnote{La loi fixait le rapport entre l'or et l'argent à 15,5, soit 4,5g d'argent pur pour 0,29g d'or pur, en continuité avec la réforme calonnienne établie en 1785.}, \oe{}uvrant de concert avec l'introduction de mesures plus rationnelles dans le régime d'émission, le pouvoir trouvait réponses à la situation hermétique sous le Directoire et à la dégradation généralisée de la monnaie\footnote{Rognage, mise en circulation de pièces affaiblies, contrefaçons et excès de monnaies de cuivre en circulation...}. De cette émanation de la souveraineté nationale découle en outre un autre fait majeur pour notre développement : la loi confirme l'abandon de la conception de monnaie de compte, d'une valeur abstraite\footnote{Découlant de sa représentation par une quantité de métal variable au gré du souverain.}, associée à la livre tournois, utilisée sous la Royauté. Monnaie de compte et monnaie réelle se confondent, et alors que l'histoire retenait le trouble des travaux de l'Assemblée constituante concernant le papier-monnaie\footnote{Les assignats créés par les décrets des 19 et 21 décembre 1789 seront sanctionnés d'un cours de monnaie à la lettre de l'article 3 des décrets du 16 et 17 avril 1790, et connaitront par la suite une large dépréciation du fait de leur usage courant par l'état à court de liquidités. Entre 1790 et 1793, l'assignat perd 60\% de sa valeur, conduisant à l'adoption d'une succession de lois restrictives, concernant notamment la fin de la publication des taux de change en 1793, et la mise en place d'un régime répressif visant à interdir les différences de traitement entre la monnaie métallique et la monnaie papier, notamment dans les opérations courantes de paiement. Dès le 8 septembre 1793, la Terreur déclarera la non-acceptation de l'assignat  passible de la peine de mort (\bibentry{white1952}), et l'inflation incontrôlée dans une période de conflit opposant la France et le gouvernement Britannique, imposera la suppression des assignats le 18 mars 1796, et leur démonétisation le 4 février 1797.}, la loi du 24 germinal an \textsc{xi} accordera pour quinze années le monopole de l'émission du papier-monnaie à une jeune société anonyme formée à Paris le 18 janvier 1800, la Banque de France. De sa qualité de personne privée découlait une multitude d'arguments en faveur de la confiance et du consensus autour des billets émis, et la comparaison naturelle pouvant être effectuée avec la structure du capital de la Banque d'Angleterre n'aboutissait à aucun rapport d'équivalence du fait que la Banque de France, \og par la nature des opérations auxquelles elle se restreint, ne court aucune chance de se trouver en avance avec le Gouvernement : elle n'aura pour débiteurs que des particuliers solvables et contraignables \fg{}\footnote{Note insérée au \emph{Moniteur} du 7 pluviôse an \textsc{viii} (27 janvier 1800).}. On observe là une première mutation organique de la monnaie importante pour la compréhension de notre développement. Le billet de banque, émis par une personne privée et convertible en or au titre d'une quantité spécifiée de métal, a perduré de manière linéaire sur l'intégralité de la période des guerres napoléoniennes\footnote{En opposition avec la Banque d'Angleterre qui suspendra la convertibilité pour préserver la stabilité des réserves d'or.} et favorisera le développement économique de la France sous le Second Empire ainsi que la modernisation du système bancaire. La monnaie, et ce qui semble confirmer notre propos antérieur, en tant que phénomène social, présenterait historiquement des éléments permettant de soutenir la thèse de son caractère détachable en partie de l'autorité étatique, malgré la conservation non sans difficultés de la nécessité de rattachement à un élément matériel imprimant de sa valeur intrinsèque\footnote{Dans le cas étudié, le métal précieux, pour la raison de la tendance logarithmique de sa découverte, abaissant sa disponibilité relative avec l'écoulement du temps et son exploitation.}.\par

\setstretch{1}\subsubsection{De l'avénement du système monométallique en Europe}\setstretch{1.42}
 
 En cette matière toutefois, le sens commun s'est rapidement trouvé oblitéré et la contrainte physique s'effacera progressivement de la pensée économico-juridique. À cet effet, le processus de transition d'un étalon-or vers la crise des changes se démarque par ses enseignements sur capacité du consensus à conférer le caractère monétaire indépendamment de toute exigence de valeur intrinsèque, et démontre l'influence des extranéités sur la structure organique des monnaies nationales. Les prémices de l'avènement du seul l'étalon-or, en lieu et place du système \emph{bimétallique}\footnote{Le terme de \emph{bimétallisme} connaitra une remise en cause en 1858 et son utilisation s'avèrera souvent abusive, de sorte qu'il se matérialisera généralement sous une forme dégénérée conduisant à un monométallisme de fait, induit par le comportement des agents économiques au regard des discordances entre le rapport légal fixe et le rapport commercial résultant des cours libres de l'or et de l'argent. (\bibentry{macleod1858}).}, trouveront en ce sens leurs origines aux états-Unis, dès 1848\footnote{\bibentry{pareto1982}.}, ou le taux de conversion or-argent, surévaluant tendanciellement l'argent par rapport à la demande d'or indispensable à la satisfaction des besoins du commerce international, entraînera une érosion de la réserve américaine et sera à l'origine d'une importante découverte d'or en Californie. En quelques années, on constatera une production mondiale multipliée par six\footnote{\bibentry{goldstats1949}.}, et la baisse du \gls{taux d'escompte} résultant de l'augmentation de l'encaisse-or, se matérialisera comme un catalyseur du crédit, conduisant, au moment de l'appauvrissement des gisements et de la déliquescence des infrastructures d'exploitation, à une crise du système bancaire international\footnote{La plus importante rareté relative de l'argent par rapport à l'or conduira à un dépassement de la valeur légale de l'argent par sa valeur physique. Les agents économiques se retrouvent dans une situation ou il devient préférable d'acheter de l'argent avec de l'or au cours légal, afin de le revendre sur les marchés, la banque centrale réglant la différence de valeur par la liquidation de ses réserves d'argent. La reconstitution des réserve entraine un mécanisme plus complexe : la banque centrale se trouve en une situation contraignante, et l'acquisition impérieuse d'argent au prix des marchés conduit les agents à spéculer sur sa réévaluation. Le gain à l'échange couvrant les coûts de refonte, de frappe et de transport, la demande de \gls{thésaurisation} des monnaies d'argent entrainera à moyen terme une crise des règlements par la diminution relative de circulation des signes monétaires. (\gls{Loi de Gresham}, 1858).}. Le 13 octobre 1857, les banques américaines suspendront les paiements en argent, et la frappe de monnaie d'argent française diminuera de près de 93,5\% entre 1956 et 1957\footnote{Données (en millions de francs) : 54,4 en 1856, 3,47 en 1857.}. La crise des règlements résonnera en Europe avec la convocation de la Convention monétaire du 20 novembre 1965 visant à harmoniser le poids des monnaies nationales et rétablir l'intercirculation des monnaies d'argent, et finalement, la Convention monétaire de 1878 abandonnera la frappe des monnaies d'argent dans l'\emph{Union Latine}\footnote{\bibentry{pareto1992}.}. Cette dérivée tendancielle vers le monométallisme de droit puisera en réalité ses origines dès 1871, à la suite guerre franco-prussienne de 1870, et de l'adoption de la \emph{Gesetz, betreffend die Ausprägung von Reichsgoldmünzen}\footnote{Traduction : Loi relative à la frappe de pièces d'or impériales.} du 4 décembre, contraignant le versement de l'indemnité de guerre due par l'état français de s'effectuer en une nouvelle monnaie : le mark-or\footnote{Le mécanisme sous-jacent permettait l'obtention d'une meilleure stabilité relative des taux de change, favorisant les rapprochements entre la livre britannique et la monnaie allemande tout en assurant l'équilibre de la balance des paiements par micro-ajustements permanents, indolores pour les populations. \bibentry{brulenaudet2021}.}. Les arguments en faveur de la fonction centralisatrice de la monnaie cristalliserons dans la politique d'Otto von Bismarck au travers de la loi 9 juillet 1873, et le mark deviendra la monnaie fédérale de l'Empire, défini par un poids d'or de 0,398g et complété par l'attribution du monopole d'émission des billets à la Reichsbank en 1875\footnote{\bibentry{berstein2014}.}.\par
 
 \setstretch{1}\subsubsection{Vers une compétence de fait attribuée à la banque centrale pour la création de valeurs monétaires}\setstretch{1.42}
 
 Toutefois, on s'aperçoit que les contestations autour du système monométallique naissent en réalité dès la décennie suivant son introduction en Europe\footnote{Il est par ailleurs surprenant d'observer que déjà les interrogations concernant les conséquences probables de la baisse des reserves d'or, émergeaient en 1857. \bibentry{chevalier1857}.}. L'idée afférente véhiculée notamment, au sein de la doctrine germanique, par Mitchell-Innes et Georg Friedrich Knapp\footnote{\bibentry{knapp1905}.}, trouvera une application pratique en 1914, avec un constat rapidement éprouvé par la Guerre : le système ne permet pas de faire correspondre la quantité de monnaie en circulation aux besoins illimités de l'économie militaire\footnote{\emph{Ibid., p.1.}}. De cette observation s'est vue sacralisée par la contrainte une nouvelle théorie de la valeur, justifiée par le principe de souveraineté entendu au sens strictement juridique, selon lequel le pouvoir pourrait librement définir l'étendue de sa propre compétence et refuser l'indexation fixe de sa monnaie sur une quantité prédéterminée d'or. On perçoit là une seconde mutation majeure de la monnaie pour notre développement. La création de valeur transcende la conception métallique afin de contrer les effets de la raréfaction du medium physique devant les besoins de circulation : la formation \emph{ex materia}, devient \emph{ex nihilo} par la planche à billets, et on observe qu'il existerait des externalités qui légitimeraient la mutation d'\gls{avilissement} de la monnaie. Cette précellence de la compétence étatique, bien qu'heureuse en cas de crises de la production, apporte néanmoins une série de questionnements au regard des problématiques de justice sociale. Pourrait-on considérer que l'impression \emph{irrationnelle}\footnote{En ce sens qu'elle se matérialise par une déconnexion de droit et de fait avec tout support métallique.} de monnaie fiduciaire puisse s'interpréter en une imposition imperceptible ? À cette interrogation, la réalité historique semble répondre par l'affirmative\footnote{\bibentry{keynes1919}.}, et ce, pour satisfaire les désidératas de l'économie belliqueuse\footnote{Mathématiquement, financer la Guerre par la contribution des agents économiques aurait nécessité des prélèvements de 50 à 60\% sur les revenus, la production ou la consommation, ce qui aurait induit mécaniquement une baisse de la propension à consentir à l'impôt, malgré la charge d'utilité publique. Finalement, il aurait été probable que cet état de tension aboutisse à une pacification, voire, une anti-militarisation d'une tranche de la population, alors que couvrir par la voie monétaire, les dépenses, se présentait comme une solution moins douloureuse de contribuer à l'effort de guerre. Preuve : en 1917, 82\% de la dépense publique trouvait fondement auprès de la defense nationale, absorbant 30 milliards de francs, alors qu'en 1913, le \acrshort{PIB} de la France était estimé à 60 milliards avec une tendance vraissembable à la baisse durant le conflit. On en déduit que le chiffre de 60\% d'imposition n'est probablement pas disproportionné. \bibentry{fontvieille1976}.}. Ce constat transparait aujourd'hui dans l'étude comparée de la comptabilité nationale des différents états parties majeures au conflit : alors que le Royaume-Uni avait favorisé le recours à la contribution individuelle, l'imposition par tête aurait connu une croissance de près de 175.95\% entre 1913 et 1918, contre 14.44\% en France, avec une situation fondée sur la paracentèse du pouvoir d'achat par l'inflation\footnote{\bibentry{clerc2013}.}. Le choix d'abandonner le système étalon-or légal pour octroyer une compétence de fait à la Banque centrale pour la création de valeurs monétaires est, nonobstant, hautement questionnable, notamment car il apparait conduire à une situation démocratiquement instable, ou semble remise en cause la réalité du constat, par le contribuable, de la nécessité de la contribution publique, de l'assiette, du recouvrement et de la durée.\par
 
 En conséquence, la mutabilité de la monnaie, lumineusement prouvée les événements historiques de l'ère contemporaine qui ont participé à la fondation de sa théorie générale, nous offre à présent deux axiomes pour son approche. D'une part, il serait possible de s'accorder sur l'existence de points de détachement avec l'entité étatique\footnote{D'après ce que nous enseigne la naissance des prérogatives de la Banque de France et de sa qualité de personne morale de droit privé.}, au profit d'une approche institutionnaliste, dont le fondement ne se situe guère dans l'État mais dans l'existence d'un lien social induit par la consécration d'une communauté de paiement. D'autre part, du fait de la fiduciarisation progressive de la monnaie\footnote{Processus mis en exergue depuis août 1914 en Europe.}, notre réflexion sur  l'interprétation phénoménologique de cette-dernière, suggérant sa potentielle virtualité n'apparait pas comme une idée dénuée de sens. Alors que l'étude technique des monnaies-cryptographiques abouti à une division franche de la doctrine autour de la qualification monétaire, cette démonstration préalable approfondie se présentait nécéssaire avant toute apprehension juridique de ces nouveaux objets. Toutefois, la lecture que nous préconisons ne s'accorde pas parfaitement avec le travail de droit produit par la jurisprudence et le législateur, que nous proposerons de remettre en perspective afin de déterminer les problématiques fiscales émergentes. À cet effet, il conviendra d'analyser ce que nous nommerons le \emph{paradoxe de la qualification}, à la lumière des écrits jurisprudentiels du droit interne et des efforts d'harmonisation au niveau communautaire.

\setstretch{1}\subsection{Une qualification trouble fruit d'un effort de réflexion sur la durée}\setstretch{1.42}

Juridiquement, reste un fait, en apparence, bien ésotérique à observer, que celui de l'absence de définition légale stricte de la monnaie\footnote{\bibentry{libchaber1992}.}. Le Code monétaire et financier se restreindra à une formule évasive : \og La monnaie de la France est l'euro \fg\footnote{Article L. 111-1 du Code monétaire et financier (CMF).}, sans pour autant émettre de suggestions précises quant aux critères de la monnaie qu'il nous a été permis d'évoquer. Néanmoins, cette ouverture théorique permet à l'inverse, une interprétation extensive de la qualification d'objet monétaire, laissée à la libre appréciation des juges et de la doctrine, qui suscite l'interrogation aujourd'hui, à l'égard les mutations organiques que connait la valeur grâce à l'avènement des nouvelles méthodes de transmission des données. Dans le cadre de notre développement, nous reviendrons nécessairement sur des aspects techniques et informatiques des monnaies-cryptographiques, apportés de manière à éclairer le juriste sur certains traits essentiels à la compréhension des enjeux afférents et la mise en place de dispositions légales pour préserver les intérêts des parties concernées par cette avancée technologique, pour une raison : la compréhension juridique des crypto-monnaies (ou cryptomonnaies) est latente, tant sur le plan jurisprudentiel que légal ou réglementaire, apportant des réponses fragiles en droit interne.\par

\setstretch{1}\subsubsection{La monnaie-cryptographique, une représentation numérique de valeur privée de la garantie d'une banque centrale}\setstretch{1.42}

Le système juridique interne emprunte de façon quasiment invariable le terme d'\emph{actif numérique}\footnote{\og Pour l'application du présent chapitre, les actifs numériques comprennent[...] \fg. Article L. 54-10-1 du Code monétaire et financier (\acrshort{CMF}).} pour qualifier au sein d'un ensemble globalisant, les monnaies-cryptographiques, ou \emph{jetons de protocole}, les jetons d'investissement (\emph{security tokens}) et les jetons ouvrant droit d'usage ou d'accès à un produit ou service (\emph{utility tokens})\footnote{Cette distinction technique est issue en partie du droit européen depuis la proposition de Règlement du Parlement Européen et du Conseil sur les marchés de crypto-actifs, modifiant la directive (UE) 2019/1937, du 24 septembre 2020, et ne connait pas d'équivalent au sein du Code monétaire et financier à ce jour.}, ainsi que \og toute représentation numérique d'une valeur qui n'est pas émise ou garantie par une banque centrale ou par une autorité publique, qui n'est pas nécessairement attachée à une monnaie ayant cours légal et qui ne possède pas le statut juridique d'une monnaie, mais qui est acceptée par des personnes physiques ou morales comme un moyen d'échange et qui peut être transférée, stockée ou échangée électroniquement \fg{}\footnote{Article L. 54-10-1, deuxième tiret du Code monétaire et financier (CMF) : \og Toute représentation numérique d'une valeur qui n'est pas émise ou garantie par une banque centrale ou par une autorité publique, qui n'est pas nécessairement attachée à une monnaie ayant cours légal et qui ne possède pas le statut juridique d'une monnaie, mais qui est acceptée par des personnes physiques ou morales comme un moyen d'échange et qui peut être transférée, stockée ou échangée électroniquement.\fg{}}.

\setstretch{1}\paragraph{Les apports de la loi relative à la croissance et la transformation des entreprises dans l'effort de qualification}\setstretch{1.42}

 L'éclaircissent définitionnel \textemdash{} dont nous nuancerons la portée dans la suite de notre développement \textemdash{} opéré par la loi PACTE relative à la croissance et la transformation des entreprises, \no 2019-486 du 22 mai 2019 nous enseigne une multitude d'informations alors essentielles pour la compréhension de la notion de \emph{jeton}, et à la lecture de l'article L. 552-2 du Code monétaire et financier, s'interprète en la sorte
\og tout bien incorporel représentant, sous forme numérique, un ou plusieurs droits pouvant être émis, inscrits, conservés ou transférés au moyen d'un dispositif d'enregistrement électronique partagé permettant d'identifier, directement ou indirectement, le propriétaire dudit bien\fg. En réalité, plusieurs observations surgissent à notre appréciation, et au prisme d'une interprétation littérale stricte de la loi, une confusion semble découler d'elle même entre les biens incorporels visés à l'article L. 552-2 et les représentations numériques de valeur acceptées en paiement, alors que la méthode cryptographique fondamentale dont il est fait l'usage dans le cadre de l'inscription au dispositif d'enregistrement électronique partagé (DEEP), présente d'innombrables similitudes. Une monnaie-cryptographiques est-elle un jeton suivant la définition donnée par  l'article L. 552-2, ou est-elle un objet réaliste doué de critères spéciaux qui renverse le processus de qualification juridique traditionnel ? On aurait espéré que les efforts produits par le législateur afin d'éclairer le domaine émergent des actifs numériques aboutissent à davantage de diaphanéité et une meilleure prise en considération de la nature polymorphique\footnote{\bibentry{couppey2019}.} et polyfonctionnelle des objets susvisés. De sa qualité incorporelle, représentée par une virtualité parfaite et une capacité entière d'identification induite par la publicité relative au sein du DEEP d'un numéro unique associé à chaque portefeuille, il devient complexe de ne pas intégrer les monnaies cryptographiques au sein de l'ensemble général des jetons. Pourtant, les textes issus du travail législatif effectuent une distinction qui semble non équivoque, entre les deux premiers syntagmes dans la structure logique de l'article L. 54-10-1\footnote{Par l'intermédiaire de son renvoi à L. 552-2 du Code monétaire et financier (CMF).}, préférant isoler les jetons, consistant en \og tout bien incorporel représentant, sous forme numérique, un ou plusieurs droits\fg{}, et les \og représentations numériques d'une valeur qui peut être transférée, stockée ou échangée électroniquement\fg{}. À la notion de \emph{bien incorporel} semble s'être substituée celle \emph{représentation numérique} d'une valeur, et éclairé par notre propos antérieur sur l'espace définitionnel de la \emph{chose} et du \emph{bien}, il devient légitime de s'interroger sur les conséquences de cette différentiation au sein de la définition de ce qui semble être l'une des rares traces légales de qualification des monnaies-cryptographies. Est-ce là une consécration de l'approche institutionnaliste de la monnaie qui semble prévaloir dans la conception législative ? Cela ne s'avère que faiblement probable. Néanmoins, il ne fait aucun doute que derrière cette définition creuse se révèle une \oe{}uvre de surface, symptôme d'une déliquescence de la loi\footnote{\bibentry{fluckiger2007}.}, et à cet égard, on se laisse surprendre par la construction terminologique dénuée de cohérence dont fait état l'écrit du législateur. Par l'énumération au sein de l'article L. 552-2 du CMF, des verbes \og émis, inscrits, conservés ou transférés\fg{}, la loi démontre une imprécision majeure attestant de son inaptitude à s'approprier méthodiquement de ces nouveaux actifs : le terme \og inscrit\fg{} absorbant \emph{a priori}, par nature et par fonction les autres membres. Cette inconsistance ne nous parait toutefois pas involontaire, et au regard du volume des travaux antérieurs à l'adoption de la loi, il serait infiniment réducteur de suggérer une simple confusion de rédaction. 

\setstretch{1}\paragraph{Pour une lecture critique du rapport Landau et de ses conclusions}\setstretch{1.42}

Cette formule témoignerait d'avantage, quoiqu'imparfaitement, de l'incapacité d'anticipation par le droit positif, des enjeux juridiques autour des monnaies-cryptographiques telle qu'elle résultait du rapport Landau au Ministre de l'Économie du 4 juillet 2018\footnote{\bibentry{landau2018}.}, et de la tentative de s'approprier de l'ensemble des situations en réponse à la complexité du système étudié, pourtant au détriment de la sécurité juridique\footnote{À ce titre, l'incohérence et la complexité excessive induites par l'énumération sans véritables fondements logiques pourraient être remises en perspective à l'égard de l'objectif de valeur constitutionnel d'accessibilité et d'intelligibilité des lois (\bibentry{conseiletat2006}).}. L'ambivalence soulevée par cette assertion cristallise en réalité les trois objectifs de politique publique visés par le rapport et atteste précisément de la tension entre perspective de développement et objectif de régulation : d'une part, l'état du droit ne permet pas une appréhension pleinement efficace des crypto-monnaies, notamment en raison de la décentralisation et du polymorphisme des protocoles, d'autre part, l'environnement se retrouve à l'interface entre l'opportunité de constitution d'un avantage comparatif dans le champ du commerce juridique international et la nécessité de circonscription des risques systémiques liés en partie à la surexposition du secteur financier à ces nouveaux objets. On peine pourtant à véritablement comprendre l'esprit du travail préparatoire tant ce dernier consacre l'\emph{ambiguïté} comme approche\footnote{\og [...] il faut sans doute accepter de vivre temporairement dans une certaine ambiguïté \fg{}. \emph{Ibid}. p.44.}, si bien qu'en différant toute tentative de qualification rigoureuse, la pluralité des risques mentionnés d'une classification jugée \og trop hâtive \fg{}\footnote{\og [...] celui de figer dans les textes une évolution rapide de la technologie ; celui de se tromper sur la nature véritable de l'objet que l'on réglemente ; celui d'orienter l'innovation vers l'évasion réglementaire\fg{}. \emph{Ibid}. p.44.} laisse oublier la tendance croissante des agents à recourir à ces nouveaux actifs et le rôle pourtant élémentaire du législateur de participer à l'amoindrissement des vides juridiques. On le perçoit, les recommandations fondées semblent éminemment polarisées par une approche en termes de compétitivité du secteur, qui résonne nettement dans le texte du CMF. En réponse à l'insertion d'un nouveau chapitre visant à parfaire la compréhension de la galaxie des jetons, on regrettera la faiblesse technique dont fait état le travail de légifération. Cette conception permet-elle d'entrevoir une mince ouverture théorique autorisant la qualification des monnaies-cryptographiques au sein d'une globalité analogue à celle des devises traditionnelles ? Rien n'est moins sûr, d'autant que le silence de la loi concernant le processus de création et d'extinction des \emph{unités}\footnote{Le terme est ici à considérer avec prudence, l'\emph{unité} visant à caractériser l'élément de base d'un ensemble de caractère structuré. Larousse. (s. d.). Unité. dans ​\emph{Dictionnaire en ligne}.​ Consulté le 6 mars 2022.} ne laisse aucunement transparaitre une véritable volonté de régulation. Néanmoins, retenant les concepts dont nous souhaitions éclairer les fondements théoriques, le droit positif n'exclut nullement l'hypothèse d'un objet privé d'attaches au système de classification juridique ancestrale, et à cet effet, l'étude du processus de qualification des \emph{données personnelles} se révèle riche en enseignements\footnote{Lecture complémentaire : \bibentry{castex2018}.}. 

\setstretch{1}\paragraph{Les enseignements du droit des données personnelles}\setstretch{1.42}

Le rapprochement entre les débats doctrinaux relatifs au droit des données personnelles et des monnaies-cryptographiques surprend tant il dépeint une ressemblance entre les fractures au sujet de la reconnaissance ou non de la qualité de \emph{bien} attachée à un objet d'une relative novation, et la solution proposée par le législateur en cette matière semble davantage apporter la lumière sur nos interrogations antérieures. Précisément, dans le cadre de l'effort de règlementation des données \emph{post mortem}, deux logiques antagoniques trouvaient à s'opposer du fait des finalités morales attachées à la transmission des données à caractère personnel : la première, prévalant originellement au sein du droit étasunien, visait, dans une dimension protectrice, à considérer ces données comme partie intégrante du patrimoine transmissible, d'où découlait une rétro-qualification de fait en un \emph{bien} entendu au sens commun ; la seconde, dont les sources se localisent majoritairement en France et dans les états de tradition austro-germanique\footnote{À ce titre, citer le droit de la Hongrie serait un exemple pertinent au regard de la tradition pénale de protection de la mémoire dont il fait état.}, suggérait inversement une volonté de préservation des intérêts du défunt en matière de vie privée, soustrayant les données personnelles de toute prétention à la transmission non consentie. À cette discussion, le principe d'\emph{informationelle Selbstbestimmung}\footnote{Traduit par équivalence en \og autodétermination informationnelle \fg{}. Droit consacré par l'arrêt du Tribunal constitutionnel fédéral de la République fédérale d'Allemagne, \og Volkszählungsurteil \fg{} du 15 décembre 1983.}, dont les prémices notionnels se distinguent à la lecture attentive de l'article 40-1 de la loi du 6 janvier 1978\footnote{Article 40-1 de la loi \no 78-17 du 6 janvier 1978 relative à l'informatique, aux fichiers et aux libertés, modifié par la loi \no 2016-1321 du 7 octobre 2016) : \og Toute personne peut définir des directives relatives à la conservation, à l'effacement et à la communication de ses données à caractère personnel après son décès. Ces directives sont générales ou particulières.\par
Les directives générales concernent l'ensemble des données à caractère personnel se rapportant à la personne concernée [...] \fg{}.
}, relative à l'informatique, aux fichiers et aux libertés, estompe les interrogations quant à la conception retenue en droit positif et ne présente à notre sens aucune ambiguïté : malgré le clivage doctrinal visant à régler l'épineuse question de la qualification des données personnelles, l'évidence semble se situer dans l'émergence d'une nouvelle abstraction, partiellement intransmissible, s'émancipant en ce sens du régime général des \emph{biens}\footnote{Alors que l'article 724 du Code civil dispose que \og les héritiers désignés par la loi sont saisis de plein droit des biens, droits et actions du défunt \fg, la jurisprudence confirmera la position du législateur adoptée dans le cadre de la loi dite \og Informatique et Liberté\fg{}. Conseil d'État, 10e chambre, 18 Novembre 2021 – \no 448729 : \og Il résulte des articles 84 et 86 de la loi du 6 janvier 1978 relative à l'informatique, aux fichiers et aux libertés, que le droit d'accès aux données personnelles s'éteint au décès de la personne concernée et que, par exception, les héritiers de la personne concernée peuvent exercer, après son décès, le droit d'accès à ces données dans la mesure nécessaire à l'organisation et au règlement de la succession du défunt, en l'absence de directives relatives à la communication des données à caractère personnel de la personne décédée, ou de mention contraire dans de telles directives\fg{}.
}. Appliquer un raisonnement parallèle en vertu de la définition de la nature et du régime des monnaies-cryptographiques, et on le comprend aisément, se présenterait probablement comme une réponse efficace afin de traiter prématurément le phénomène de concurrence des monnaies dont nous explorerons les fondements et évaluerons la probabilité d'émergence, pourtant, cette hypothèse ne semble que partiellement appréhendée par la jurisprudence.\par 

\setstretch{1}\subsection{Une classification jurisprudentielle finalement analogue à celle de la \emph{monnaie légale} à la lecture croisée de la doctrine administrative et civile}\setstretch{1.42}

\og Quoi que l'on fasse, les lois positives ne sauraient jamais entièrement remplacer l'usage de la raison naturelle dans les affaires de la vie. Les besoins de la société sont si variés, la communication des hommes est si active, leurs intérêts sont si multipliés, et leurs rapports si étendus, qu'il est impossible au législateur de pourvoir à tout\fg{}.\footnote{Discours préliminaire du premier projet de Code civil, présenté en l'an \textsc{ix} par J.-E.-M. Portalis, F.-D. Tronchet, F.-J.-J. Bigot-Préameneu et J. de Maleville, membres de la commission nommée par le Gouvernement.}\par

À cet effet, le droit des monnaies cryptographiques interpelle par l'infime volume de décisions questionnant la problématique de qualification et les solutions juridiques formulées par le dialogue des juges. En droit interne, deux arrêts retiendront à ce titre notre attention pour les éléments généraux qu'ils disposent : la décision du Conseil d'État, 8\up{ème} et 3\up{ème} chambres réunies, du 26 avril 2018, \no 417809, et celle du Tribunal de commerce de Nanterre du 26 février 2020, \no 2018F00466. Il est intéressant d'observer à quel point la doctrine administrative a exprimé relativement tôt sa volonté d'appréhender ces nouveaux crypto-actifs au prisme d'une approche fiscale. 

\setstretch{1}\subsubsection{Une réponse trouble du juge administratif à l'égard de la définition des monnaies-cryptographiques}\setstretch{1.42}

À dater de 2014, des essais de définition se trouvent consacrés au Bulletin officiel des finances publiques (\acrshort{BOFiP}) et seront cités par le juge administratif ; les apports à la connaissance juridique marqueront le droit embryonnaire avec toutefois la mise en évidence d'incohérences résiduelles qui persisteront au fil de la réflexion. Ainsi, le paragraphe \no 1080 des commentaires administratifs publiés le 11 juillet 2014 au BOFiP - \emph{impôts} sous la référence BOI-BNC-CHAMP-10-10-20-40, repris en des termes identiques dans le cadre de l'actualisation du 3 février 2016 \textemdash{} quand bien même ce dernier s'attacherait à un objectif certain de clarté \textemdash{} présente au sein de ses deux premiers ensembles syntaxiques une confusion dont la persistance surprend tellement les idées contradictoires véhiculées se succèdent avec proximité\footnote{Paragraphe \no 1080 des commentaires administratifs publiés le 11 juillet 2014 au BOFiP - \emph{impôts} sous la référence BOI-BNC-CHAMP-10-10-20-40 : \og Le bitcoin est une unité de compte virtuelle stockée sur un support électronique permettant à une communauté d'utilisateurs d'échanger entre eux des biens et services sans recourir à une monnaie ayant cours légal. Les bitcoins sont acquis soit gratuitement en contrepartie d'une participation au fonctionnement du système, soit à titre onéreux sur des plates-formes internet créées afin de permettre l'achat et la vente de bitcoins contre de la monnaie ayant cours légal. L'émission du nombre de bitcoins étant limitée et déterminée, leur acquisition en vue de leur revente procède d'une intention spéculative [...]\fg{}.}. Dans un premier temps, le texte semble s'attacher à une qualification fonctionnelle des monnaies-cryptographiques\footnote{On perçoit là parfaitement la nécessité d'une approche pluridisciplinaire du sujet intégrant le droit et les sciences économiques, dans une perspective de compréhension précise des concepts techniques.}, estimant que \og le bitcoin est une unité de compte virtuelle [...]  permettant à une communauté d'utilisateurs d'échanger entre eux des biens et services \fg. Déjà survient une interrogation : une unité de compte entendue au sens strictement littéral permet-elle de façon autonome, l'échange de biens et services ? On le comprend aisément, il procéderait davantage d'un moyen standardisé accordant la faculté de quantifier la valeur d'un élément du patrimoine, que celui d'un vecteur du pouvoir libératoire. Quoi que l'on puisse suggérer de la qualité rédactionnelle du texte, ce dernier semble cependant effleurer une dualité fonctionnelle partiellement intégrée dans l'objet d'étude, et en ce sens s'attacherait efficacement à une définition de la monnaie telle qu'elle résulte de nos développements antérieurs. En dépit de cette observation, on laissera remarquer avec une particulière attention l'usage de la notion de \emph{communauté}, dont désormais l'on ne contestera point sa substance dans la qualification monétaire. Dans un second temps, et ce qui semble témoigner d'une véritable incertitude dans l'appréhension juridique des cryptomonnaies, est affirmé que : \og les bitcoins sont acquis soit gratuitement en contrepartie d'une participation au fonctionnement du système, soit à titre onéreux sur des plates-formes internet [...] \fg, laissant la lecture en confrontation avec une formulation oxymorique dont l'incohérence peine à s'expliquer tant elle interpelle. Alors que l'ordonnance \no 2016-131 du 10 février 2016 rappelait que dans la sphère du droit général des obligations, le contrat à titre gratuit se déduit \og lorsque l'une des parties procure à l'autre un avantage sans attendre ni recevoir de contrepartie \fg{}\footnote{Article 1107 du Code civil : \og Le contrat est à titre onéreux lorsque chacune des parties reçoit de l'autre un avantage en contrepartie de celui qu'elle procure. Il est à titre gratuit lorsque l'une des parties procure à l'autre un avantage sans attendre ni recevoir de contrepartie\fg{}.}, le juge administratif se limitant à un rappel du commentaire, trouble infiniment la compréhension de la pourtant quasiment unique évocation du processus de création des unités virtuelles. À la lettre du texte, on méconnait tout du mécanisme d'inscription initiale au \og système\fg{}, dont l'on déduit que le terme précède l'actuel \og DEEP\fg{} consacré pour la première fois au sein de l'article L. 223-12 issu de l'ordonnance \no 2016-520 du 28 avril 2016\footnote{Article L. 223-12 du CMF issu de l'ordonnance \no 2016-520 du 28 avril 2016 : \og Sans préjudice des dispositions de l'article L. 223-4, l'émission et la cession de minibons peuvent également être inscrites dans un dispositif d'enregistrement électronique partagé permettant l'authentification de ces opérations, dans des conditions, notamment de sécurité, définies par décret en Conseil d'Etat\fg.}, ainsi que de la \emph{contrepartie} attendue, de sa nature et de sa proportion. À cette dernière question, la chambre commerciale de la Cour de cassation précisera d'ailleurs à l'occasion d'une décision rendue au 16 décembre 2014, \no 13-25.765, le critère du contrat à titre gratuit, par une interprétation stricte en parfaite antinomie avec la notion de contrats commutatifs \og dans lesquels les obligations du débiteur\fg{} excéderaient \og notablement celles de l'autre partie \fg{}. Est-il possible d'entrevoir un rapport d'obligations liant d'une part le sujet de droit participant \textemdash{} en une manière totalement ignorée par l'administration \textemdash{} au fonctionnement d'un \emph{système} privé de qualification juridique ? Alors que la cour semble s'attacher à une logique purement binaire, s'interroger sur la contrepartie évoquée par l'administration à l'issue d'une étude mathématique sur le processus de participation au système Bitcoin \textemdash{} reprenant la vision hautement polarisée par la jurisprudence \textemdash{} semblerait permettre, à notre sens, de compendieusement appréhender les acteurs en présence contribuant à l'inscription des unités sur le DEEP, avec l'espoir de clarifier la qualification d'un \emph{a priori} non-statut juridique.\par

\setstretch{1}\paragraph{La nécessité d'une \emph{contrepartie} au prisme d'une analyse du processus de participation au système Bitcoin}\setstretch{1.42}

Apprécier l'essence de la \emph{contrepartie} en se refusant à un travail de modélisation attesterait d'une profonde insuffisance de rigueur tant le sujet repose sur un raisonnement algorithmique propre. Nous reviendrons postérieurement sur les implications du processus démocratique mis en \oe{}uvre pour la vérification des transactions effectuées sur le réseaux, mais dans un premier temps, analysons brièvement le paradigme associé à la production de la masse monétaire. À la lecture du \emph{livre blanc} \og Bitcoin: A Peer-to-Peer Electronic Cash System\fg{}\footnote{\bibentry{nakamoto2008}.} dépeignant les fondements théoriques d'une solution au problème de \gls{double dépense}\footnote{Le problème de double dépense renvoie, en informatique des systèmes décentralisés, à une situation ou un même débiteur serait en position de signer en envoyer sans intervalle de temps, deux transactions prenant pour point d'entrée une unique unité de monnaie-cryptographique, à destination de deux bénéficiaires distincts. En conséquence, la probabilité pour que deux ensembles de \emph{n} n\oe{}uds connaissent d'une information divergente sur l'état de la chaîne de transactions pendante accentue le risque de désaccord sur la validité du registre à l'échelle du réseau.}, plusieurs éléments retiennent notre attention : premièrement, toute institution décentralisée, pour satisfaire à son objectif d'autoriser l'échange de jetons entre des usagers, se retrouve contrainte à la mise en fonction d'un service de fourniture de \emph{preuves de date}\footnote{\bibentry{gorkhali2020}.} reposant sur l'information publique d'empreintes numériques\footnote{Notées \emph{Hash} au sein du livre blanc et dans la pratique informatique.} associées à un \emph{bloc} de transactions. Chaque \emph{bloc} reprenant l'empreinte de son précédent, ainsi qu'une valeur signée des éléments qui le composent, une modification même infime de la date d'enregistrement d'une transaction ou de la valeur de celle-ci ne devient envisageable que si un consensus autorise la mutation de tous les \emph{blocs} postérieurs. Sous une version simplifiée faisant l'omission volontaire de certaines variables considérées comme non essentielles dans le cadre de notre développement, un \emph{bloc} minimal pourrait s'illustrer ainsi :

\begin{lstlisting}
bloc {
	index: 250,
	timestamp: "2022-03-10T15:56:23.798",
	previousHash: "9e74c712ec2a60946eC16be2ccf05fac627823d14537e78c049a7114344b0309",
	hash: "95e7030a5fb80fa14c4c9f10252ef71f3408b83955b390e1a23e64077194aec5",
	data: [liste des transactions]
}
\end{lstlisting}
On le perçoit parfaitement, la compréhension de la structure générale d'un \emph{bloc} n'est pas si complexe et on identifie clairement au sein de notre modèle les informations essentielles au fonctionnement du registre : la position relative du \emph{bloc} au sein de la chaine\footnote{Position que l'on qualifiera d'\emph{index}, permettant notamment de pérenniser une relation de précédent au sein de la chaîne et de traiter les conflits en cas de diffusion simultanée de plusieurs blocs valides.}, sa date de validation, son emprunte numérique et celle de son précédent, avec pour suite une liste de transactions dont la longueur sera décidée arbitrairement ou algorithmiquement sur décision du développeur du protocole. Deuxièmement, il nous devient possible de produire avec précision une chronologie visant à isoler la phase d'émission des unités représentatives de valeur et le fait générateur de l'acquisition originaire. Ainsi, chaque détenteur du client pair à pair\footnote{Un \gls{système pair à pair} (\emph{peer-to-peer}) décrit un modèle d'échange en réseau où chaque entité est à la fois client et serveur. Dans le cadre du Bitcoin, ce système peut être qualifié de parfaitement décentralisé du fait que les connexions entre participants se réalisent sans intervention d'une infrastructure de répartition (tiers de confiance). Lecture complémentaire : \bibentry{fox2001}.} d'accès à la chaîne représentant un n\oe{}ud, l'algorithme connaissant des caractéristiques techniques de la machine sur lequel il est exécuté, collecte et agrège un nombre prédéterminé de transactions, vérifie leur validité et procède à l'élaboration d'une \og preuve de travail\fg{} (\emph{Proof of Work})\footnote{Le terme de \og preuve de travail\fg{} n'est toutefois pas général et tend à une décélération dans les usages (du fait notamment de la croissance exponentielle de la consommation électrique attachée au développement du réseau), au profit de certains protocoles reposant sur de nouveaux paradigmes, à la manière de la \og preuve d'enjeux\fg{} (\emph{Proof of Stake}). Toutefois, on estime qu'en octobre 2016, 90\% de la capitalisation boursière totale des monnaies-cryptographiques existantes dépendait de la preuve de travail. \bibentry{gervais2016}.} prenant comme arguments d'entrée, l'emprunte du bloc précédent, l'emprunte du bloc à traiter et une variable traditionnellement nommée le \emph{nonce}, seul élément d'ajustement visant à la satisfaction des conditions de validation posées par le protocole. Là se situe la contrepartie requise au participant : le protocole fixe une valeur \emph{cible} représentée sur 256 bits\footnote{Soit un nombre parmi 2\up{256} possibilités, la valeur maximal étant $\approx 1.1579.10^{77}$.}, la validation d'un bloc par un n\oe{}ud étant subordonnée à la recherche et à la diffusion d'une emprunte inférieure ou égale à cette valeur\footnote{En comparant récursivement une suite de variations du \emph{nonce} concaténé à l'emprunte du bloc, l'obtention d'un \emph{hash} inférieur ou égal à la cible procède d'un mécanisme pseudo-aléatoire complexifiant le processus itératif et préservant l'asymétrie du coût de calcul. À la suite d'une étude sur la minimisation de chaînes en Python 3.10 avec une fonction de hachage SHA-256, visant à aboutir à un préfixe égal à \og 000000 \fg{} sur la chaîne \og Louis Brulé Naudet \fg{}, les résultats démontrent le caractère exponentiel de la puissance de calcul nécessaire à la validation d'un bloc : alors qu'il ne nécessitera que 605 itérations afin d'aboutir à un préfix égal à \og 000 \fg{}, cette valeur sera de 6 776 916 dans le cas d'une suite de \og 000000 \fg{}.}. On le déduit aisément, cette recherche quasi-aléatoire du fait des propriétés des fonctions de hachages mises en \oe{}vre couplées à la galaxie des nombres envisageables requiert une puissance de calcul considérable issue de l'allocation par chacun des n\oe{}uds d'une unité centrale de traitement. Envisager une durée de plusieurs années dans le cadre de la découverte par un unique participant, d'une preuve de travail valide n'apparait donc pas tant déconnecté de la réalité\footnote{\bibentry{dagannaud2017}.}, si bien que la fondation Bitcoin préconise l'utilisation d'une configuration matérielle minimale\footnote{Au 14 mars 2022, la configuration recommandée était, pour un ordinateur exploité par une version récente de Windows, MacOS ou Linux : 7 gigabytes d'espace libre sur le disque de stockage, accessibles à une vitesse de lecture et d'écriture d'au moins 100 Mb/s, couplés à deux gigabytes de mémoire RAM, une connexion internet autorisant une vitesse de téléversement d'au moins 50 kilobytes par seconde et une durée de mise à disposition minimale de 6 heures par jour. Étant précisé que le téléchargement d'usage représente un volume de données équivalent à environ 20 gigabytes par mois, après un téléchargement de 340 gigabytes lors de la configuration du n\oe{}ud.} afin d'optimiser la durée et la consommation énergétique demandées pour l'inscription d'un bloc validé sur le registre. Se fait alors jour la nature de la participation au réseau et cette digression technique ne fait que souligner l'évidence : l'attribution ne suit en rien une contribution à titre gratuit comme évoquée par l'administration. Pourtant, l'idée d'une monnaie énergétique surprend tant elle fait état de réminiscences dans la pensée économique, et déjà en 1921 les interrogations sur ses fondements théoriques trouveront à être documentées\footnote{L'idée était alors de considérer qu'une certaine quantité d'énergie produite en une heure équivalait à un dollar américain, et pouvait être librement transférée entre les états et les civilisations, avec pour finalité la substitution parfaite au système du monométallisme-or supposé catalyseur des conflits armés. \bibentry{ford1921}.}. Finalement, les n\oe{}uds exprimeront leur acceptation du bloc diffusé par le participant en travaillant à générer le suivant, usant de l'empreinte numérique associée à ce premier, et, dès l'instant précédant son ajout au registre, le participant par l'intermédiaire de son adresse\footnote{Adresse : représentation cryptographique sur 160 bits de la clé publique associée au portefeuille d'actifs numériques du participant.} et des conventions du protocole, réalisera algorithmiquement l'insertion à l'indice 0 du bloc, d'une transaction privée d'émetteur dont il est le destinataire\footnote{\og By convention, the first transaction in a block is a special transaction that starts a new coin owned by the creator of the block\fg. \emph{Ibid.} p.4.}. Cette inscription\footnote{Préférons le terme d'\emph{inscription} à celui de \emph{transaction} afin d'éviter toute confusion avec les dispositions de l'article 2044 du Code civil.} fondamentalement particulière révèle alors le fait générateur de la création monétaire décentralisée, et il devient difficile de comprendre les raisons pour lesquelles tant la jurisprudence, que la loi et la doctrine ignorent ce mécanisme pourtant porteur de toutes les solutions relatives à la compréhension du processus d'augmentation de la masse monétaire d'une devise cryptographique.

\setstretch{1}\paragraph{L'absence d'émission organisée par une personne morale}\setstretch{1.42}

 S'il est inscrit sur un support électronique et doué d'un caractère par nature négociable\footnote{Dans le cadre d'une acquisition dérivée à la suite d'un transfert de propriété entre deux utilisateurs du réseau.}, le Bitcoin n'est indéniablement pas le fruit d'une émission organisée par une personne morale mais résulterait davantage du procédé démocratique et de l'autorité de fait accordée à l'algorithme de validation : aucun droit de créance ne saurait naître de cette situation. Le doute concernant l'existence d'un lien d'obligations s'en trouve débordé et rien ne semble s'opposer à la relation d'équivalence qui existerait entre le langage et la monnaie telle que nous l'évoquions. La reconnaissance non contraignante par un système conventionné au sein d'un groupe génère sans ambages une situation où chacun se retrouve en mesure de quantifier, comparer et s'échanger un élément incorporel, dont il serait méconnaître du fonctionnement même des systèmes, d'en considérer l'absence de valeur intrinsèque. Cette théorie entre alors en collision avec toute proposition d'assimilation des monnaies-cryptographiques à la catégorie juridique des instruments financiers ou des monnaies électroniques. Les uns étant par nature limités au critère organique rattaché à l'émission étatique ou privée par l'intermédiaire d'une personne morale, d'un fonds commun de placement, de financement ou de titrisation\footnote{Article L.211-2 du CMF : \og Les titres financiers, qui comprennent les valeurs mobilières au sens du deuxième alinéa de l'article L.228-1 du code de commerce, ne peuvent être émis que par l'État, une personne morale, un fonds commun de placement, un fonds de placement immobilier, un fonds professionnel de placement immobilier, un fonds de financement spécialisé, ou un fonds commun de titrisation\fg.}, les autres représentant une créance sur un émetteur contre remise de fonds à la lettre de l'article L.315-1 du CMF\footnote{Article L.315-1 du CMF : \og  La monnaie électronique est une valeur monétaire qui est stockée sous une forme électronique, y compris magnétique, représentant une créance sur l'émetteur, qui est émise contre la remise de fonds aux fins d'opérations de paiement définies à l'article L. 133-3 et qui est acceptée par une personne physique ou morale autre que l'émetteur de monnaie électronique\fg{}.}.\par
 
 \setstretch{1}\paragraph{Comprendre le mécanisme lié à l'offre de monnaie de long terme}\setstretch{1.42}
 
 Un dernier point de tension semble également émerger de la lecture du commentaire administratif concernant l'émission des unités sur un équivalent décentralisé de la notion de \emph{marché primaire} et cristalliserait injustement la conception d'équilibre de long terme entre l'offre et la demande de monnaie-cryptographique. À cet effet, dans son deuxième considérant, le juge rappelant que \og l'émission du nombre de bitcoins étant limitée et déterminée, leur acquisition en vue de leur revente procède d'une intention spéculative\fg{} opacifie davantage ce qu'était l'état de la compréhension des devises cryptographiques. Certes, cette assertion est justifiable par l'analyse du mécanisme de création inhérent au protocole Bitcoin, la relativité du principe impose toutefois plusieurs limites dont une brève analyse permet de souligner l'insuffisance dans une perspective d'appréhension réelle. La logique mathématique responsable de la première inscription au registre d'une unité nouvellement revendiquée par un n\oe{}ud est fascinante tant elle sous-tend un concept bio-mimétique complexe pourtant mis en \oe{}uvre d'une manière remarquablement accessible : afin de se conformer avec la représentation logarithmique découlant de la production d'or à l'échelle de l'écoulement du temps, le protocole a, dès sa conception, adopté un \emph{politique monétaire} modélisée par une suite géométrique de raison $\frac{1}{2}$ convergeant vers une masse monétaire totale à l'échéance théorique de maturation du protocole équivalant  à 20 999 999,977 BTC. Précisément, le schéma suit une tendance cyclique marquée par une transition (\emph{\gls{halving}}) tous les 210 000 blocs\footnote{La difficulté étant ajustée afin d'aboutir à la validation d'un bloc tous les intervalles de dix minutes environ, notamment pour des raisons de diminution du temps de traitement des transactions sur le réseau, chaque cycle possède une durée approximative de quatre années. Lecture complémentaire : \bibentry{grunspan2020}.}, divisant le coefficient de production par 2, avec au temps 0 une valeur de 50 unités revendiquées pour chaque bloc validé par le participant, telle que l'offre maximale $\varphi$ (\emph{\gls{maximum supply}}) se détermine par la suivante\footnote{L'offre maximale ici représentée s'exprime en une sous-division de la devise Bitcoin, le Satoshi dont le rapport de grandeur est 1 SAT = $1.10^{-8}$ BTC.} :
 
\begin{equation}
	\varphi=\frac{\sum_{i=1}^{32}210000[\frac{50*10^8}{2^i}]}{10^8}
\end{equation}
\pgfplotsset{ every non boxed x axis/.append style={x axis line style=-},
     every non boxed y axis/.append style={y axis line style=-}}
\pgfplotsset{compat=1.8}
\begin{center}
\begin{tikzpicture}
\begin{axis}[
    axis lines = left,
    grid style=dashed,
    width=11cm, height=6cm,
    scaled y ticks = false,
    y tick label style={/pgf/number format/sci},
    xlabel = {$blocs$},
    ylabel = {$Bitcoin$},
]
\addplot [
    ]
    coordinates {
    (52500, 2625000)(105000, 5250000)(157500, 7875000)(210000, 10500000)(262500, 11812500)(315000, 13125000)(367500, 14437500)(420000, 15750000)(472500, 16406250)(525000, 17062500)(577500, 17718750)(630000, 18375000)(682500, 18703125)(735000, 19031250)(787500, 19359375)(840000, 19687500)(1050000, 20343750)(1260000, 20671875)(1470000, 20835937)(1680000, 20917968)(1890000, 20958984)(2100000, 20979492)(2310000, 20989746)(2520000, 20994873)(2730000, 20997436)(2940000, 20998718)(3150000, 20999359)(3360000, 20999679)(3570000, 20999839)(3780000, 20999919)(3990000, 20999959)(4200000, 20999979)(4410000, 20999989)(4620000, 20999994)(4830000, 20999997)(5040000, 20999998)(5250000, 20999999.35950000)(5460000, 20999999.67240000)(5670000, 20999999.82780000)(5880000, 20999999.90550000)(6090000, 20999999.94330000)(6300000, 20999999.96220000)(6510000, 20999999.97060000)(6720000, 20999999.97480000)(6930000, 20999999.97690000)
    };
\end{axis}
\begin{axis}[
      width=11cm, height=6.5cm,
      axis y line*=right,
      axis x line*=top,
      xlabel = {$Période$},
      ylabel = {$\phi$},
      x axis line style={draw opacity=0},
      y axis line style={draw opacity=0},
      every axis label/.append style ={opacity=0},
      ]
\addplot[
	dotted
    ]
    coordinates {
    (2009, 1)(2012, 2)(2016, 3)(2020, 4)(2024, 5)(2028, 6)(2032, 7)(2036, 8)(2040, 9)(2044, 10)(2048, 11)(2052, 12)(2056, 13)(2060, 14)(2064, 15)(2068, 16)(2072, 17)(2076, 18)(2080, 19)(2084, 20)(2088, 21)(2092, 22)(2096, 23)(2100, 24)(2104, 25)(2108, 26)(2112, 27)(2116, 28)(2120, 29)(2124, 30)(2128, 31)(2132, 32)(2136, 33)(2140, 34)
	};
\end{axis}
\end{tikzpicture}
\end{center}

Analytiquement, on retrouve le seuil de déclenchement du \emph{halving} fixé à 210 000, le montant initial de la première inscription au registre ainsi que le coefficient de progression géométrique. Une observation légitime pourrait alors questionner les effets d'une variation de ce seuil sur l'offre maximale de Bitcoin et à ce titre, peu important le montant initialement implémenté au sein du protocole, la relation avec l'offre est telle qu'imaginer un système ayant pour ambition une offre immensément plus importante est parfaitement envisageable. Pour exemple \emph{in vivo}, l'Ether (\acrshort{ETH}) ne dispose pas d'une offre maximale sur le long terme mais simplement un plafond annuel et la pluralité des protocoles en développement interroge quant à l'applicabilité du raisonnement effectué par l'administration. Quelles seraient les conséquences de l'émergence d'une monnaie-cryptographique structurellement stable ? À la lumière de la jurisprudence administrative, il reste impossible d'en établir les contours. En dépit d'un contenu très embarrassé, les fondements théoriques posés par le Conseil d'État, témoignent néanmoins d'un premier effort de sacralisation des problématiques liées aux monnaies cryptographiques dans la pensée juridique française, et s'il est vain de considérer que cette décision fonde efficacement le critère de ces nouveaux objets, on ne peut nier son influence sur la doctrine judiciaire.\par

 \setstretch{1}\subsubsection{La monnaie-cryptographique, un bien \emph{fongible} et \emph{consomptible} à la manière de la monnaie légale}\setstretch{1.42}

À ce titre, la décision du Tribunal de commerce de Nanterre du 26 février 2020, \no 2018F00466 précise davantage d'éléments en vue d'une qualification véritablement pratique des crypto-monnaies, et cette décision marque ce qui semble être la première immixtion majeure de la question devant l'ordre judiciaire. En l'espèce, alors que la chaîne de blocs \emph{Bitcoin} faisait l'objet d'une scission\footnote{Le \og \gls{hard fork} \fg{} en bases de données décentralisées consiste en un changement du paradigme de validation d'ensembles de transactions, partagé par tous les noeuds du réseau, pouvant conduire à une division permanente de la chaîne de blocs en cas de concurrence entre différentes versions logicielles.} (\og hard fork\fg\footnote{\bibentry{swan2015}.}) donnant naissance au Bitcoin Cash (ci-après \acrshort{BCC}),  nouvelle crypto-monnaie indépendante, revendiquée automatiquement par l'intermédiaire du protocole aux détenteurs de Bitcoin (ci-après \acrshort{BTC}) à date du 1\up{er} août 2017, la volonté de requalification des conventions en contrats de prêt à usage\footnote{Article 1875 C. civ. : \og Le prêt à usage est un contrat par lequel l'une des parties livre une chose à l'autre pour s'en servir, à la charge par le preneur de la rendre après s'en être servi\fg{}.} et le soutien d'un enrichissement injustifié ont précocement émergé de la vie des affaires au vu de l'opportunité représentée par l'acquisition originaire de cette nouvelle crypto-devise. Mais aux fins de connaitre de la nature du contrat de prêt de BTC, il convenait de consciencieusement appréhender l'\oe{}uvre de qualification de l'objet de la convention, dans une décision riche de référentiels et points de réflexion. 

 \setstretch{1}\paragraph{Le Bitcoin, un bien consomptible car aliénable: démonstration et limites}\setstretch{1.42}

En premier lieu, et après un rappel de l'article 1892 du Code civil\footnote{Article 1892 C. civ. : \og Le prêt de consommation est un contrat par lequel l'une des parties livre à l'autre une certaine quantité de choses qui se consomment par l'usage, à la charge par cette dernière de lui en rendre autant de même espèce et qualité\fg{}.}, le juge, consacrant le caractère consomptible du Bitcoin \og lors de son utilisation, que ce soit pour payer des biens ou des services, pour l'échanger contre des devises ou pour le prêter, tout comme la monnaie légale, quand bien même il n'en est pas une \fg{}, soulève en réalité une série de questionnements en dépit d'une solution dont on ne pourrait se méprendre sur son efficacité. À lire la décision, on retrouve la complexité rattachée à la monnaie par notion de \emph{consomptibilité} qui interpelle tant elle se recherche dans les insuffisances d'une approche matérielle découlant d'une distinction incidente au sein du Code civil\footnote{Les termes de biens \emph{consomptibles} et \emph{non consomptibles} ne font l'objet d'aucune référence explicite au sein du Code civil et à cet aune, déjà le droit romain abordait la notion de \emph{resquae primo usuconsumuntur} afin de décrire les choses dont on ne pouvait faire usage sans les consommer. \bibentry{lambertye2020} (27 septembre 2020).}. Alors que la consommation par l'usage de liqueurs ne semble guère porter à confusion\footnote{Article 587 C. civ. : \og Si l'usufruit comprend des choses dont on ne peut faire usage sans les consommer, comme l'argent, les grains, les liqueurs, l'usufruitier a le droit de s'en servir, mais à la charge de rendre, à la fin de l'usufruit, soit des choses de même quantité et qualité soit leur valeur estimée à la date de la restitution\fg{}.}, penser celle de la monnaie, à la lettre de l'article 587 du Code civil requiert la manipulation de bien des abstractions conceptuelles, dont le raisonnement ne laisse s'effacer les contradictions. Si la décision de la première chambre civile de la cour d'appel de Besançon du 3 novembre 2015, \no 14/00743\footnote{Confirmée par l'arrêt de la deuxième et quatrième chambre réunies de la cour d'appel d'Aix-en-Provence du 6 févr. 2019, \no 16/17584.} regardait la monnaie comme un bien consomptible par nature du fait que l'on ne peut, en principe, en faire l'usage que par l'aliénation\footnote{\bibentry{libchaber1992}.}, on retrouve une traduction juridique formelle de notre analyse d'une monnaie dont on avait vidé le critère symbolique : privée de la notion d'aliénation la faisant disparaitre du contenu du patrimoine, elle ne représente plus aucune source d'utilité pour son propriétaire ou son détenteur du fait de son impossible transmission entre les personnes effaçant elle-même la nécessité de reconnaissance mutuelle autour de l'objet monétaire. La monnaie deviendrait alors inductivement un objet douée d'une consomptibilité juridique\footnote{\bibentry{bonfils2003}.} et non plus matérielle\footnote{On parle parfois également de consomptibilité naturelle.}. Cette \emph{consomptibilité subjective}\footnote{\bibentry{mercier2005}.} en ce sens qu'elle découlerait de la volonté de qualifier en tant que tel un objet d'étude en vue de l'appréhension que l'on en fait, suppose pourtant de multiples caractéristiques, avec en premier lieu, une considération essentielle à l'égard de sa valeur pécuniaire\footnote{\bibentry{jaubert1945}.}. Ainsi, supposer que la consomptibilité juridique apporterait à la monnaie la présomption d'une qualité de réserve de valeur\footnote{Dont la reconnaissance présente un important passé historique en doctrine, et permettrait de rejoindre la pensée économique.} participe à un raisonnement sous-jacent particulièrement sensible aux interrogations. Focaliser sa réflexion sur l'aliénation en tant que source de la consomptibilité amène à réfléchir sur la définition même du référentiel dans lequel on se place. S'il ne fait désormais plus de doute sur la possibilité d'aliénation d'un immeuble, peut-on également lui imprimer la qualité de bien consomptible ? En la matière toutefois, il semblerait que le consensus ne laisse de place au débat tellement la non-consomptibilité de l'immeuble semble consacrée en droit civil\footnote{\bibentry{simler2018}}. Or, serait étrangement trouble, la pensée qui viserait à affirmer que l'immeuble est à la fois consomptible et non-consomptible\footnote{\bibentry{leclerc2009}.}. Est-ce à penser que la consomptibilité induite par l'aliénation serait l'apanage de la monnaie ? Cette conclusion, et on le comprend aisément, serait la plus heureuse pour toute réflexion visant à consacrer l'analogie entre les monnaies cryptographiques et les monnaies traditionnelles. Néanmoins, cela pourrait également s'entendre en la simple affirmation de la fragilité du raisonnement initial visant à considérer que la consomptibilité, quand bien même elle peut s'entendre en un sens subjectif, naît de la possibilité d'aliénation. Une chose reste pour le moins certaine : la monnaie détient juridiquement des particularités qui lui sont propres et essentielles à sa qualification.\par

\setstretch{1}\paragraph{L'acceptation d'un pouvoir libératoire en tant qu'instrument de paiement}\setstretch{1.42}

Évoquer la consomptibilité du Bitcoin d'une manière équivalente dans les fondements théoriques, à celle de la monnaie, ressert également les liens entre monnaie et instrument de paiement. Bien qu'il soit admis de manière générale que juridiquement, le paiement ne représente que l'exécution d'une obligation restreinte dans sa portée par l'effet relatif du contrat\footnote{\bibentry{carbonnier1988}.}, pouvant s'effectuer discrétionnairement entre les parties, l'analogie introduite par la décision alors qu'elle caractérise la consomptibilité du BTC par son usage amène naturellement à s'interroger sur la terminologie employée au regard de la confusion qui peut résulter de son interprétation : est faite mention de la possibilité de \og payer des biens ou des services [...] tout comme la monnaie légale \fg{}, alors que s'en tenir à une définition générale du paiement, voire, suggérer un autre exemple d'instrument utilisé pour cet usage, aurait été sans doute plus judicieux, surtout afin d'éviter que la lecture entre en contradiction avec l'article 1343-3\footnote{Article 1343-3 C. civ. : Le paiement, en France, d'une obligation de somme d'argent s'effectue en euros.} du Code civil. Si l'on s'accorde sur le fait qu'un moyen de paiement\footnote{En d'autres termes, d'extinction d'obligation.} non monétaire puisse exister de manière indépendante, le critère de consomptibilité, tel qu'il est présenté, à savoir, de manière exclusive pour la réalisation de paiements, prêts ou échanges, met en exergue une identité avec la définition donnée par le Doyen Carbonnier\footnote{\bibentry{carbonnier1995}.}, précisant que l'on \og ne peut utiliser des instruments monétaires qu'en les dépensant, en les impliquant dans un paiement\fg{}. Alors que l'on pourrait être amené a limiter la qualification du BTC en un moyen de paiement entendu au sens du Code monétaire et financier\footnote{Article L311-3 CMF : Sont considérés comme moyens de paiement tous les instruments qui permettent à toute personne de transférer des fonds, quel que soit le support ou le procédé technique utilisé.} possédant un pouvoir libératoire même partiel\footnote{Conditionné à l'acceptation du créancier dans le cadre de la relation contractuelle.} cette solution ne fait qu'enrichir l'argumentaire autour d'une consécration du BTC comme une monnaie autonome, et déjà en 2018, le rapport Landau se plaçait à l'interface entre la pensée économique et juridique, se faisant point de convergence des critiques\footnote{\bibentry{bonneau2018}.}.\par

 Malgré les difficultés de qualification dont témoigne ce dernier, une de ses grandes qualités résiderait dans sa considération des monnaies-cryptographiques en tant que \emph{monnaies privée}, alors appartenant à une conception séparée du soutien d'une quelconque banque centrale visant à l'impression d'un cours légal. De nombreux auteurs ont alors argumenté de la non pertinence d'une considération de pouvoir libératoire, en ce sens que leur utilisation est nécessairement restreinte par l'usage des créanciers acceptant le libellé d'une opération monétaire en monnaie-cryptographique. On affirmait anciennement de son pouvoir libératoire qu'il était alors bien moins légal que conventionnel\footnote{\bibentry{corbioncondé2014}.}, néanmoins, en généralisant le caractère conventionnel à une population plus importante d'usagers, on retrouve un schéma qui a existé en droit français avant l'introduction de l'Euro. Plus précisément, la question se posait en 1994, de savoir un paiement en écu était marqué d'un caractère libératoire ou si sa fonction se restreignait à celle d'une unité de compte. À cela, la réponse ministérielle\footnote{\bibentry{assemblee1994}.} s'est montrée parfaitement claire, précisant que malgré l'existence d'une distinction entre les deux types de monnaies, et la seule reconnaissance d'un pouvoir libératoire légal au franc, l'écu étant qualifié de devise au regard de la règlementation des changes, \og il devient possible \---{} sans aucune restriction \---{}  de stipuler expressément en écu, en matière de droit des obligations, et les clauses d'un contrat prévoyant le règlement dans cette devise entre deux résidents ne peuvent être frappées de nullité d'ordre public\fg{}\footnote{\emph{Ibid. p.2602}}. Bien que l'extension législative ne concernait alors que l'écu, et se trouvait limitée par l'impossibilité pour les administrations de percevoir le recouvrement de recettes publiques en une devise autre que légale, cette illustration témoigne de l'influence d'un usage généralisé d'une monnaie prévue contractuellement sur sa qualification. Il ne faut toutefois pas se méprendre, l'usage ne créer la monnaie que de manière relative, et à ce titre, il est intéressant de s'interroger sur l'évolutivité du processus de qualification, ou la variable temporelle ne saurait jouer un rôle minime. Certes, en 2014, le pouvoir libératoire d'une monnaie-cryptographique se trouvait en grande partie limité par sa non reconnaissance partagée entre les acteurs de l'économie, mais cette vision trouve aujourd'hui à être nuancée, d'autant plus que si l'on s'attache à une observation historique de l'obligation monétaire, l'ancien droit faisait usage d'une diversité de monnaies de compte et de paiement\footnote{\bibentry{carbonnier2004}.}, dont la tendance contemporaine tend à agréger les fonctions\footnote{\bibentry{libchaber1992}.}. Ainsi, on cherchera systématiquement le critère de l'unité de valeur\footnote{\bibentry{lassalas1997}.}, additivement à la fonction d'instrument de paiement, dans une perspective de qualification d'une chose monétaire. Sur le plan théorique, cette capacité prend source dans l'exigence de fongibilité afin de conserver la qualité de support nécessaire à l'accumulation et au transfert de la valeur, en particulier du crédit par la qualité de vecteur des intérêts.\par 

\setstretch{1}\paragraph{Le respect de l'exigence de fongibilité : critère déterminant pour la conception d'une unité de compte théorique}\setstretch{1.42}

Dans sa décision, le tribunal répond à une problématique largement partagée en doctrine concernant la fongibilité de la monnaie cryptographique. Alors que certains auteurs soutiennent la non équivalence des BTC entre eux du fait d'une individualisation\footnote{\bibentry{rousille2015}.} implémentée dans le protocole lors de la validation des transactions, le juge avance au contraire, qu'ils font \og l'objet d'un rapport d'équivalence [..] permettant d'effectuer un paiement au sens où l'entend l'article 1291 ancien du Code civil, devenu l'article 1347-1\footnote{Article 1347-1 C. civ. : \og Sont fongibles les obligations de somme d'argent, même en différentes devises, pourvu qu'elles soient convertibles, ou celles qui ont pour objet une quantité de choses de même genre\fg{}.}\fg{}. En réalité, on peut même parfaitement remettre en cause les théories qui prônaient l'individualisation par le protocole informatique du fait du fonctionnement du DEEP, traitant simplement de manière équivalente l'ensemble des unités de monnaie cryptographique inscrites. Cette qualification n'est toutefois pas sans effet pour détermination de la nature des cryptomonnaies. En consacrant l'équivalence réciproque, cela ouvre la conception théorique d'une unité de compte au sens ou il sera possible de procéder à une évaluation quantitative de la valeur des choses du commerce. L'unité devient divisible et interchangeable, soutenue par un système de numérotation utilisant la base 10 jusqu'à 8 décimales\footnote{Pour le seul cas du Bitcoin, ou 1 satoshi = \(1.10^{-8}\) BTC.}, dont la similitude avec l'article L. 111-1 du Code monétaire et financier\footnote{Article L.111-1 du CMF : \og La monnaie de la France est l'euro. Un euro est divisé en cent centimes\fg{}.} est marquante dans sa logique. Cette analogie est d'ailleurs celle retenue en droit interne par la doctrine fiscale, qui analysera le Bitcoin comme \og une unité de compte virtuelle stockée sur un support électronique permettant à une communauté d'utilisateurs d'échanger entre eux des biens et services sans recourir à une monnaie ayant cours légal \fg{}\footnote{BOI-BNC-CHAMP-10-10-20-40-20160203, no 1080, 3 févr. 2016.}. Mais cette considération ne semble pas emporter un parfait consensus, et dans une publication de 2018\footnote{\bibentry{bdf2018}.}, la Banque de France émettra d'ailleurs un raisonnement qui lui est propre et qui marquera la forte dépendance entre les sphères économique et juridique, en élevant au rang d'axiome que la volatilité de leurs cours ne peut conduire à la caractérisation d'une unité de compte et d'une réserve de valeur. Ce raisonnement laisse également place à une interprétation plus critique, toujours au prisme de l'analyse temporelle, et on notera que ce dernier critère est en mesure de sembler injustement utilisé, notamment par l'observation analogue de valeurs mobilières dont le cours est sensible par définition aux fluctuations, qui restent pourtant qualifiées de réserves de valeur.\par

En somme, et après une nécessaire remise en respective du critère fonctionnel de la monnaie, on ne peut que remarquer une galaxie de similitudes dont la pertinence tend à s'accroitre avec l'activité économique des individus qui font société. La monnaie cryptographique semble partager, après bien des discussions en doctrine et dans la pratique des institutions, les critères propres à la monnaie légale, notamment sa consomptibilité et sa fongibilité. Toutefois, certaines caractéristiques des monnaies cryptographiques persistent en proie à la nécessité d'un effort d'interprétation, notamment concernant la relativité du pouvoir libératoire leur étant attaché, et en ce sens, il devient fondamental d'alimenter notre analyse par une étude de la crise de confiance qui touchait les systèmes du DEEP, ainsi que son évolution dans le temps afin de tirer des conclusions dans un terme prospectif. À cette fin, le second temps de notre raisonnement vise à interpréter les conséquences de ces similitudes dans le cadre d'une concurrence des monnaies entre crypto-devises et monnaies traditionnelles, notamment en vue de l'aboutissement du processus de stabilité attaché à un type encore en perfectionnement de monnaies cryptographiques : les stablecoins, en particuliers algorithmiques, ainsi que de l'influence du droit européen et du droit fiscal pour la conduite d'une politique monétaire conciliante des intérêts antagonistes d'attractivité des investissements internationaux et de souveraineté monétaire.
	
\setstretch{1}\chapter{Concurrence et lutte de monopole : entre qualification implicite et opportunités de développement, quels enjeux à la lumière d'une fiscalité incitative ?}\setstretch{1.42}

En référence aux travaux de l'école ordolibéraliste de l'Université de Fribourg-en-Brisgau\footnote{Lecture complémentaire : \bibentry{commun2019}.}, visant lumineusement à la réunion entre juristes et économistes afin d'affirmer les jalons de l'ordre économique dans sa dimension juridique, le raisonnement adopté dans ce second développement aura vocation à expliciter les rapports transversaux entre cadre juridique et paradigme économique, à l'aune de l'influence fiscale et réglementaire sur les usages de monnaies cryptographiques et \emph{in fine} le degré de concurrence entre les devises. Intitulé \emph{concurrence et lutte de monopole} au miroir de l'ouvrage \emph{Wettbewerb und Monopolkampf} de Franz Böhm\footnote{\bibentry{bohm1933}.}, il aspire d'une part, à étudier les opportunités et risques de développement d'une concurrence des monnaies sous l'angle de la pensée de l'école autrichienne, dans la perspective d'un perfectionnement des protocoles informatiques permettant un meilleur contrôle du phénomène inflationniste et de l'effritement de la crise de confiance qui touchait les monnaies cryptographiques, en confrontation avec le droit positif consacrant la monnaie en tant qu'élément de souveraineté étatique (2.1). Et d'autre part, il interrogera l'empreinte du droit fiscal interne et européen dans l'effort de qualification (2.2) afin d'aboutir à un ensemble de suggestions en fonction des diverses appréciations de politiques monétaires laissées aux institutions.\par

Évoquer la concurrence des monnaies au prisme d'une lecture croisée de l'école de Fribourg et des travaux de Friedrich Hayek se présente en un travail riche d'enseignements pour de multiples raisons. Premièrement, en renforçant la prégnance de la raison scientifique dans l'érection d'une constitution économique se faisant l'interdépendance entre les politiques monétaires, fiscales et commerciales\footnote{On retrouve majoritairement là l'exaltation de la réflexion euckienne telle qu'elle ressort de ses \emph{Die Grundlagen der Nationalökonomie} de 1940}, on s'écarte du nominalisme juridique selon lequel le droit s'efforce de graviter autour de l'individu sans considération de l'harmonie entre les hommes faisant société\footnote{\bibentry{villey2013}.}. Il devient alors réalisable l'étude plus approfondie des phénomènes concurrentiels, les ordonnateurs\footnote{Entendus en toutes institutions douées d'un pouvoir d'influence sur les marchés, y compris sur la politique monétaire et sa stabilité.} cristallisant l'ossature juridique assurant le juste fonctionnement des marchés. Alors que Walter Eucken admirait le \emph{wettbewerbsodrnungen}\footnote{Wettbewerbsodrnungen (allemand) : ordre concurrentiel.} pour son fondement finaliste visant à la conclusion des actions économiques quotidiennes entre les agents, on ne peut que trouver là une justification claire de ce que l'on pourrait traduire en une émergence de moyens permettant une fluidification du commerce juridique entre les personnes. On observe en ce sens un croisement entre la logique poursuivie par l'école de l'ordolibéralisme, qui trouvera sa concrétisation en matière monétaire en 1976 avec la publication de \emph{Pour une vraie concurrence des monnaies} de Friedrich Hayek : le degré de pouvoir d'un monopoleur ne peut satisfaire au bien-être de la nation, quelle que soit sa forme\footnote{En ce sens, Wilhelm Röpke écrira notamment au sein de son \emph{Die Lehre von der Wirtschaft} de 1963 (traduit par \og doctrine de l'économie\fg{}) : \og Le fait décisif reste que les monopoleurs disposent d’un degré de pouvoir sur leurs marchés et sur l’économie, et qu’un système économique bien ordonné basé sur une juste relation entre la performance et la récompense ne peut pas tolérer \fg{}. \bibentry{fevre2017}.}. Penser le réalisme d'une concurrence des monnaies n'est d'ailleurs pas un fait historiquement dénué de fondements, et nous évoquerons (2.1.1.2) en la matière les opportunités théoriques qu'il devient possible de dessiner à la lumière de la statique du droit positif.

\setstretch{1}\section{Perfectionnement des protocoles et conception moderne de la souveraineté étatique : un antagonisme renforcé au profit de l'accroissement du degré concurrentiel entre les devises}\setstretch{1.42}

Il ne faut toutefois pas se méprendre sur les termes employés et on ne peut se limiter à une étude sous l'angle d'un \emph{droit de la concurrence} tel qu'il résulte de la loi, du fait de l'absence de mécanismes juridiques positifs d'appréciation d'un écosystème pluri-devises\footnote{Certes, l'article L.511-4 du CMF prévoit l'inapplicabilité des dispositions relatives aux pratiques anti-concurrentielles exercées par des établissements de crédit, il ne faut toutefois pas là entendre à une sur-interprétation concernant la production de monnaie, dont on sait qu'elle est légalement convenue en tant que monopole étatique. \og Les articles L. 420-1 à L. 420-4 du code de commerce s'appliquent aux établissements de crédit et aux sociétés de financement pour leurs opérations de banque et leurs opérations connexes définies à l'article L. 311-2, aux établissements de monnaie électronique pour l'émission et la gestion de monnaie électronique et leurs opérations mentionnées à l'article L. 526-2 ainsi qu'aux établissements de paiement pour leurs services de paiement et leurs services connexes définis à l'article L. 522-2\fg{}.}. Néanmoins, il serait également réducteur de nier l'existence conceptuelle d'une rivalité dans les usages, dans la mesure où, placée au sein d'une approche conçue comme exclusivement fonctionnelle, l'offre d'\emph{objets} en qualité d'intermédiaires des échanges et d'unités de compte peut présenter une diversité selon la reconnaissance d'un certain nombre d'individus imprimant un caractère libératoire. Ainsi, alors que l'on reprochait aux monnaies cryptographiques, et ce, notamment au sein du rapport Landau\footnote{\og Les monnaies sont aussi des réserves de valeur. Cette qualité est souvent déniée aux crypto-monnaies, car leurs cours sont très volatils : en moyenne vingt-cinq fois plus que les actions américaines, cinq fois plus que les matières premières et douze fois plus que la devise japonaise, le yen. Il est évident que les crypto monnaies ont nourri, au cours de la dernière année, une forte activité spéculative\fg{}. \emph{Ibid. p.12.}}, cette poussée tendancielle à la volatilité, qu'en serait la lecture des juridictions et institutions dans la perspective d'une atteinte à maturité des protocoles permettant de s'affranchir des contraintes liées aux règles majoritairement fixes de l'inscription originaire au registre permettant la structuration de la base monétaire ? On le comprend aisément, et cette notion structure en réalité notre thèse, cela se présenterait comme la condition minimale afin d'envisager la naissance d'une reconnaissance collective portant en elle les germes de la solution au problème qualificationnel. Pourtant, une seconde approche semble conduire la réflexion en droit français, dont le fondement se retrouve consacré par la décision du 26 février 2020 du tribunal de commerce de Nanterre, vidant quasiment de sa substance tout effort de créativité en la matière concernant l'esquisse d'un espace monétaire dual

\setstretch{1}\subsection{Le paradoxe de la qualification : un rejet de l'horizontalisation des rapports monétaires entre les personnes}\setstretch{1.42}

À la lecture de la décision, un terme particulier guide le raisonnement du juge dans sa qualification, et le conduit à appliquer une distinction spécifique entre le Bitcoin et la monnaie que l'on pourrait qualifier de \og légale \fg{}\footnote{\og Que le BTC est \og consommé \fg{} lors de son utilisation, que ce soit pour payer des biens ou des services, pour l'échanger contre des devises ou pour le prêter, tout comme la monnaie légale, quand bien même il n'en est pas une\fg{}. Nanterre, 6e ch., 26 févr. 2020, \no 2018F00466.}. S'interroger sur cet adjectif permet d'observer la singularité de la solution : le juge avance la sacralisation de l'approche organique en organisant une rupture claire avec l'évaluation fonctionnelle au profit de la compétence régalienne.\par

\setstretch{1}\subsubsection{La sacralisation de l'approche organique : une rupture claire avec l'évaluation fonctionnelle au profit de la compétence régalienne}\setstretch{1.42}

Dans une vision politique générale, on comprend la motivation à considérer que la monnaie ne peut être qu'une compétence régalienne\footnote{\bibentry{bonneau2018} ; \bibentry{bonneau2019}.}. D'une part, cette conception permet de légitimer l'action de l'état de battre monnaie, d'autre part, elle élève au rang d'institution souveraine l'entité en charge de sa production. Il convient alors d'observer les fondements du rejet de toute concurrence en matière monétaire, ainsi que les conséquences de cette négation de la réalité économique pesant notamment sur l'ordre public.\par

L'opposition entre le Bitcoin et la \og monnaie légale \fg{} s'entend en réalité autour d'un élément, pierre angulaire d'autres concepts sous-jacents concernant notamment le pouvoir libératoire : la confiance des agents en l'instrument monétaire. Considérer que la monnaie se présente comme un objet standardisé assimilé à un moyen d'échange au sein d'une communauté, dont tous les membres sont d'accord concernant son utilisation\footnote{\bibentry{lietaer2004}.}, l'institutionnalise en une composante essentielle du lien social\footnote{\bibentry{dollo2018}.}. Pour ce faire, l'État possède l'appareil contraignant qu'est le Code pénal\footnote{Historiquement, une parfaite illustration était observable au sein de l'article unique de la loi du 12 février 1916, disposant que : \og En temps de guerre, toute personne convaincue d'avoir acheté, vendu ou cédé, d'avoir tenté ou proposé d'acheter, de vendre ou de céder des espèces et monnaies nationales, à un prix dépassant leur valeur légale, ou moyennant une prime quelconque, sera condamnée à une peine de six jours à six mois d'emprisonnement et à une amende de cent francs à cinq mille francs ou à l'une de ces deux peines seulement.\fg{}}, notamment en sanctionnant le refus de recevoir des pièces de monnaie ou des billets de banque ayant cours légal en France\footnote{Article R.642-3 C. pén. : Le fait de refuser de recevoir des pièces de monnaie ou des billets de banque ayant cours légal en France selon la valeur pour laquelle ils ont cours est puni de l'amende prévue pour les contraventions de la 2e classe.}. On retrouve un rôle fondamental en la monnaie qui permet la \og polarisation mimétique des efforts et des violences \fg{}\footnote{\bibentry{dosquet2018}.}, et dont la reconnaissance de la valeur forme l'identité commune au sein d'un ordre social harmonisé et stable. L'idée d'une prégnance de l'État dans la politique monétaire que nous évoquions notamment au sein de notre digression sur l'effritement du modèle étalon-or (2.2.1.2) persistera dans la pensée juridique grâce à son aptitude quant à la justification des politiques monétaires expansionnistes\footnote{Le Gouvernement pouvant ainsi stimuler la production de monnaie par l'intermédiaire de la Banque centrale en abaissant artificiellement les taux directeurs avec pour objectif de stimuler le crédit. Lecture complémentaire : \bibentry{plihon2013}.} et dans cette abstraction, la monnaie ne tient son caractère légal que par le fait de sa régulation en raison des moyens dont disposent les banques centrales envers les établissements de crédit. Par cette décision, le juge rappelle alors que le critère de la monnaie dépend intimement de l'institution qui la garantie, et sous-tend dans cette approche organique que le BTC en l'état de l'art, présente un défaut de confiance\footnote{\bibentry{kleiner2019}.} encore dirimant pour sa qualification en tant que monnaie.\par
Finalement, on retrouve assez clairement cette conception historique de la souveraineté que soutenait implicitement Jean Bodin dans son expression de \og puissance absolue et perpétuelle d'une république \fg{}\footnote{\bibentry{blanc2006}.}, le droit de battre monnaie devant appartenir à l'État en tant que marque de sa souveraineté, élevé à un niveau équivalent à celui de la loi du prince. L'étude de la pensée de Bodin se révèle en ce sens riche en éléments de compréhension pour l'analyse de la prérogative monétaire telle qu'elle est sacralisée par le juge, et permettra de consacrer théoriquement la dimension contractuelle du monnayage en tant que garde-fou du souverain\footnote{\og Il ne faut donc pas confondre la loy et le contract : car la loy depend de celuy qui a la souveraineté, qui peut obliger tous ses subjects, et ne peut s’y obliger soy même : et la convention est mutuelle entre le prince et les subjects, qui oblige les deux parties reciproquement : et ne peut l’une des parties y contrevenir au prejudice, et sans le consentement de l’autre et le prince en ce cas n’a rien par dessus le subject\fg{}. \emph{Ibid. \S{} 41.}}. En l'échange de la justice naissant du devoir de \og maintenir par la force des armes et des lois ses sujets en sûreté de leurs personnes, biens, et familles \fg{} garantie par le prince, les citoyens se trouvent soumis à la sujétion et à la foi. On aperçoit là la nécessité d'un entretien de la foi publique par l'exigence d'une certaine fiabilité dans le cours légal et la question du titre. Le c\oe{}ur de la solution de Jean Bodin est par ailleurs remarquable de sa prise en compte des éléments de l'histoire récente à ses contemporains en vue de la suggestion d'une éventuelle sanction au non respect des préoccupations morales et de la justice. En dénonçant les pratiques de Philippe le Bel (1268-1314) ayant conduit à un avilissement considérable de la monnaie\footnote{Le concours entre l'insuffisance des droits seigneuriaux et la baisse du bénéfice royal issu du droit de frappe en raison de l'augmentation du cours des métaux conduiront le roi à dévaluer sa monnaie. \bibentry{gauvard2013}.} et une crise économique et financière renforcée par la guerre contre la Flandre, il intronise la notion de \emph{roi faux-monnayeur} se devant d'être limitée par un ensemble de propositions, notamment concernant la centralisation de l'opération de monnayage et l'abolition du recours à la frappe au marteau au bénéfice de la méthode du moulage afin d'accroître la qualité et complexifier la reproductibilité des pièces. Néanmoins, s'attacher à cette conception de l'approche institutionnelle de la monnaie semble emporter quelques difficultés, avec en premier lieu l'encrage profondément métalliste du système proposé par Bodin ne répondant désormais en aucun point avec la réalité de la pratique moderne. Pourtant, un parallèle semble, à notre sens, être autorisé au regard de la centralisation opérée par le processus de construction européenne.\par

Si la cour permanente de justice internationale, dans l'affaire concernant le paiement de diverses emprunts serbe et brésilien\footnote{En l'espèce, la question portait sur le fait de savoir si les porteurs de titres d'emprunts émis par le Brésil et la Serbie, antérieurement à première Guerre Mondiale devaient être payés, non pas en francs français tels que dévalués par la dernière loi monétaire française du 25 juin 1928, mais en francs-or tels que définis par la loi du 7 germinal an \textsc{xi}. \bibentry{permanent1929affaire}} émis en France, en date du 12 juillet 1929, consacrait l'axiome de la souveraineté monétaire en droit international par la formule \og c'est un principe généralement reconnu que tout État a le droit de déterminer lui-même ses monnaies\fg{}, cette dernière constituera la pensée originaire à la lumière de laquelle se dessinera la doctrine en particulier à la lettre des travaux de Jean Carbonnier et de Dominique Carreau. Ainsi, au prisme d'une approche supra-nationale, s'affirmera la proposition selon laquelle bien que la souveraineté monétaire se trouve illimitée en ordre interne, le droit international peut y apporter des tempéraments\footnote{\bibentry{blanc2002}.}.

\setstretch{1}\paragraph{Une consécration non sans remise en cause de la dimension institutionnaliste de la monnaie}\setstretch{1.422}

Pour autant, il faut se garder de donner aux idées davantage de portée qu'elles n'en ont, et il ne semble que peu justifié d'observer la \emph{souveraineté monétaire} dans une conception absolue. De fait, afin de minutieusement appréhender les conséquences d'une parallélisation des usages monétaires, l'écueil principal se retrouverait dans sa sacralisation en tant que droit inaliénable\footnote{En conséquence, cela revient à ignorer les effets d'un abandon de la monnaie nationale au profit d'une monnaie étrangère dans le cadre d'un processus d'adoption visant à accroitre les investissements internationaux (à la manière de la dollarisation de l'Équateur de 2001), ou de la construction d'une union monétaire régionale.} sans faire état de la distinction entre le principe et ce qui relève de son exercice. Alors qu'aux termes du préambule de la constitution de 1946, \og sous réserve de réciprocité, la France consent aux limitations de souveraineté nécessaires à l'organisation et au maintien de la paix \fg{}\footnote{\bibentry{vaulpane2017}.}, on comprend qu'une lecture compendieuse tiendrait probablement pour vraie la proposition selon laquelle la souveraineté monétaire demeure, en réponse aux mécanismes de rupture prévus au sein des traités internationaux. Or, il semblerait qu'une conception plus complexe de la \emph{souveraineté monétaire} puisse émerger de l'étude réelle de la compétence, en ce sens que l'on pourrait distinguer une construction à double niveaux faisant état d'une part des prérogatives touchant les pratiques internes de compte au sens de l'établissement d'un cours légal\footnote{\emph{Ibid. p. 5.}}, de seigneuriage et d'établissement d'une symbolique monétaire, mais également une caractérisation secondaire, polarisée par la capacité à imposer les instruments monétaires au sein des pratiques internes de compte et de paiement, à concevoir et exécuter une politique monétaire et de change\footnote{Notamment concernant la capacité à influencer la valorisation externe de la monnaie interne.}, et à démontrer la maîtrise des pratiques de conversion. En procédant de la sorte, on autorise une lecture multiple de la souveraineté monétaire, notamment par l'intermédiaire des confrontations aux contraintes réglementaires de fait et de droit accompagnant les accords internationaux. Ainsi, il convient de s'interroger sur les effets supposés de la distinction, notamment en ce qui concerne la manifestation du principe d'\emph{exclusivité monétaire nationale} au prisme de la concurrence des devises.\par

La naissance du concept d'\emph{exclusivité monétaire nationale} fait paradoxalement état d'une relative proximité dans l'histoire juridique\footnote{En ce sens, Herrenschmidt [1998] illustrera la comparaison au miroir d'une \og parenthèse dans l'histoire de la monnaie\fg{}. \emph{Ibid. p. 3.}}, se rattachant pour partie à l'émergence des États-nations unifiés Italien et Allemand de 1861 et 1870\footnote{Lecture complémentaire: \bibentry{pecout2004}.}. néamoins, son apport pour la compréhension de la pratique monétaire contemporaine ne pourrait se voir limité au regard des critères qu'il apporte et auxquels il adosse un certain nombres d'objectifs attachés à la conception moderne de la souveraineté. Réceptacle d'ambitions de développement industriel et militaire, l'État-nation, contraint à l'homogénéisation des pratiques commerciales et leur rationalisation, notamment en vue de la satisfaction des exigences de l'économie belliqueuse, se fera acteur de la suppression des reliquats du féodalisme en imposant une exclusivité monétaire nationale reposant sur trois principes : l'unicité, l'exclusivité et la spécificité\footnote{\emph{Ibid. p. 17.}}. On observera dans le sillage de cette mouvance un processus d'appariement nation-monnaie et on retrouve assez aisément les fondements d'une conception moderne de la souveraineté telle qu'elle transparait dans la décision du tribunal de commerce de Nanterre : la monnaie tient son caractère de l'institution chargée de son émission.\par

Pour autant, il semblerait que ce paradigme soit confronté à une pluralité de limites, et ce, dès le lendemain de la dissolution du système de Bretton-Woods\footnote{Il sera alors remarquable de constater que le principe de l'\emph{exclusivité monétaire nationale} connaîtra sa généralisation à la presque totalité des États-nations de la planète, attestant dès lors du succès de la doctrine et de son influence sur les diverses conceptions de la souveraineté.}. Alors que le concept d'\emph{exclusivité monétaire nationale} tenait à une forme de relation de domination de la monnaie nationale sur les monnaies parallèles\footnote{\bibentry{moulier2016}.} \---{} colorant ainsi l'esprit d'une véritable concurrence des devises en considération de la politique d'évincement de tout autre objet monétaire \---{} apparaitront les germes d'une revalorisation du degré concurrentiel considéré jusqu'à lors comme quasiment nul. 





\newpage

\setstretch{1}\subsection{La crise de confiance envers les Dispositifs d'Enregistrement Electronique Partagés à t'elle vocation à perdurer ?}\setstretch{1.42}

\setstretch{1}\subsubsection{Le constat d'une inflation des usages, renforcée en 2022}\setstretch{1.42}

\setstretch{1}\subsubsection{La multiplication des modes d'acquisition}\setstretch{1.42}

\setstretch{1}\subsection{L'émergence des \emph{electronic money tokens} algorithmiques, pierre angulaire d'un système intégralement décentralisé}\setstretch{1.42}

\setstretch{1}\subsubsection{Définitions et enjeux}\setstretch{1.42}

\setstretch{1}\subsubsection{Vers la généralisation d'un facteur de risques systémiques}\setstretch{1.42}








\setstretch{1}\section{Pour une approche particulière des monnaies cryptographiques : quelle fiscalité et ses conséquences dans une approche finaliste de la concurrence des monnaies ?}\setstretch{1.42}

\setstretch{1}\subsection{La considération progressive d'une précellence européenne pour le traitement des questions crypto-monétaires}\setstretch{1.42}

\setstretch{1}\subsubsection{Éléments de droit comparé européen}\setstretch{1.42}

\setstretch{1}\paragraph{Des prémices de la régulation au Royaume-Uni, une parfaite compréhension des enjeux ?}\setstretch{1.42}

\setstretch{1}\paragraph{L'Europe des dissonances : un droit encore en attente d'unification au niveau communautaire}\setstretch{1.42}

Einkommensteuergesetz section 23

\setstretch{1}\subsubsection{L'apport novateur des instances de l'Union européenne dans la perspective de réponse à la réalité économique et fiscale}\setstretch{1.42}

\setstretch{1}\paragraph{Une assimilation au régime des opérations bancaire et financière en matière de taxe sur la valeur ajoutée au profit de la compétitives des échanges}\setstretch{1.42}

\setstretch{1}\paragraph{Vers l'adoption d'un règlement européen sur les marchés de crypto-actif (MiCA)}\setstretch{1.42}








\setstretch{1}\subsection{Opportunité ou technologie de rupture : l'influence de la fiscalité sur le développement des monnaies cryptographiques}\setstretch{1.42}

\setstretch{1}\subsubsection{L'épineuse question de la responsabilité déclarative, inhibiteur ou catalyseur des usages ?}\setstretch{1.42}

\setstretch{1}\paragraph{Les difficultés contemporaines d'appréciation des régimes et de leurs effets}\setstretch{1.42}

\setstretch{1}\paragraph{Pour nouvelle responsabilité des prestataire de services sur actifs numériques à l'égard de l'administration fiscale et ses conséquences}\setstretch{1.42}

\setstretch{1}\paragraph{Vers une omniscience de l'administration sur les réseaux publics ?}\setstretch{1.42}

\setstretch{1}\subparagraph{La propicité des fonctions de hachage et de la résistance aux collisions pour une déclaration automatisée}\setstretch{1.42}

\setstretch{1}\subsubsection{Pour une fiscalité plus incitative visant à l'alimentation de l'économie réelle}\setstretch{1.42}

\setstretch{1}\paragraph{Faciliter les dépenses courantes au profit des usagers non professionnels}\setstretch{1.42}

\setstretch{1}\paragraph{Permettre aux sociétés un usage plus opérant dans le cadre de leurs activités}\setstretch{1.42}

 Vers une soutenance des projets systémique en faveur aux monnaies nationales
 
 Les USA se reposent sur les cryptos des gafas pour soutenir le dollar, le tout grace aux marchés financiers et la banque centrale par le rachat des dettes
 
 Avant le 19e, chaque baron local battait monnaie, faire une reference historique
 
 "electronic money token" 
 
 Le bitcoin en tant que monnaie étrangère

 Propos conclusifs
\newpage

% Bibliography
\bibliographystyle{plainnat}
\bibliography{bibliographie}

% Acronyms
\printglossary[type=\acronymtype]
% Glossary
\printglossary

\end{document}
